\documentclass[11pt,a4paper,openany]{book}
\usepackage{fontspec}
\usepackage{xcolor}
\usepackage{tcolorbox}
\tcbuselibrary{skins,breakable}
\usepackage{booktabs}
\usepackage{longtable}
\usepackage{titlesec}
\usepackage{fancyhdr}
\usepackage{geometry}
\usepackage{setspace}
\usepackage{polyglossia}
\usepackage{hyperref}
\setmainfont{Noto Sans}
\newfontfamily\arabicfont[Script=Arabic]{Noto Naskh Arabic}
\setmainlanguage{english}
\setotherlanguage{arabic}
\let\textAR\textarabic
\definecolor{bgpage}{HTML}{FAF6F0}
\definecolor{bgcard}{HTML}{F5EFE3}
\definecolor{bgquran}{HTML}{EDE5D8}
\definecolor{bghadith}{HTML}{E8EDE5}
\definecolor{primary}{HTML}{1B4D3E}
\definecolor{primarydk}{HTML}{10332B}
\definecolor{accent}{HTML}{C9A857}
\definecolor{accent2}{HTML}{B8964A}
\definecolor{textmain}{HTML}{1C1917}
\definecolor{textmuted}{HTML}{57534E}
\geometry{top=2.5cm, bottom=2.5cm, left=2.2cm, right=2.2cm}
\pagecolor{bgpage}
\color{textmain}
\hypersetup{colorlinks=true, linkcolor=primary, urlcolor=primary, citecolor=primary, pdftitle={Surah Al-Anfal 8:24-37}, pdfauthor={Tarbiyyah Research}}
\titleformat{\chapter}[display]{\normalfont\huge\bfseries\color{primary}}{\chaptertitlename\ \thechapter}{10pt}{\Huge}
\titlespacing*{\chapter}{0pt}{0pt}{20pt}
\titleformat{\section}{\normalfont\Large\bfseries\color{primary}}{\thesection}{1em}{}
\pagestyle{fancy}
\fancyhf{}
\fancyhead[L]{\small\color{textmuted}\itshape Surah Al-Anfal 8:24--37}
\fancyhead[R]{\small\color{textmuted}\thepage}
\fancyfoot[C]{\small\color{textmuted} Tarbiyyah Research --- Quran Wiki Corpus}
\renewcommand{\headrulewidth}{0.5pt}
\newtcolorbox{quranbox}[1][]{breakable, colback=bgquran, colframe=accent, coltitle=primary, fonttitle=\bfseries\small, title={#1}, boxrule=0.8pt, arc=3pt, left=10pt, right=10pt, top=8pt, bottom=8pt, before skip=10pt, after skip=10pt}
\newtcolorbox{tafsirbox}[1][]{breakable, colback=bgcard, colframe=primary, coltitle=primary, fonttitle=\bfseries\small, title={#1}, boxrule=0.8pt, arc=3pt, left=10pt, right=10pt, top=8pt, bottom=8pt, before skip=8pt, after skip=8pt}
\newtcolorbox{muyassarbox}[1][]{breakable, colback=bghadith, colframe=primarydk, coltitle=primarydk, fonttitle=\bfseries\small, title={#1}, boxrule=0.5pt, arc=3pt, left=10pt, right=10pt, top=6pt, bottom=6pt, before skip=6pt, after skip=6pt}
\newtcolorbox{themebox}[1][]{breakable, colback=bgcard, colframe=accent2, coltitle=accent2, fonttitle=\bfseries\small, title={#1}, boxrule=0.5pt, arc=3pt, left=10pt, right=10pt, top=6pt, bottom=6pt, before skip=8pt, after skip=8pt}
%% Quran verse - use \textAR inside minipage (NOT \begin{Arabic} which breaks in tcolorbox)
\newcommand{\quranverse}[1]{{\centering\Large\textAR{#1}\par}}
\setstretch{1.15}
\begin{document}

\begin{titlepage}
\thispagestyle{empty}
\begin{center}
\vspace*{3cm}
{\color{accent}\rule{\textwidth}{2pt}}
\vspace{1.5cm}
{\Huge\bfseries\color{primary} Surah Al-Anfal}\\[0.5cm]
{\LARGE\color{primarydk} Verses 24--37}\\[1cm]
{\Large\color{accent2}\itshape A Thematic Study of\\[0.3cm] The Call to Faith, Divine Favour, and the Fate of Disbelievers}\\[1.5cm]
{\color{accent}\rule{0.6\textwidth}{1pt}}\\[0.8cm]
{\large\color{textmuted} Detailed Notes with Arabic Text,\\ Multiple Translations, Tafsir Ibn Kathir, and Tafsir Muyassar}\\[0.5cm]
{\normalsize\color{textmuted}\itshape Based on the verified Quran Wiki corpus}\\[2cm]
{\small\color{textmuted}\textbf{Sources:}\\ Arabic: Uthmani (Al Quran Cloud + Quran.com)\\ Translations: Sahih Int, Pickthall, Yusuf Ali, Asad, Basmeih\\ Tafsir: Ibn Kathir, Muyassar}
\vfill
{\color{accent}\rule{\textwidth}{2pt}}
\vspace{0.5cm}
{\small\color{textmuted} Tarbiyyah Research \quad|\quad \today}
\end{center}
\end{titlepage}

\tableofcontents
\newpage

\chapter{Introduction}

\section{Scope and Methodology}
This study presents a detailed thematic analysis of \textbf{Surah Al-Anfal (8), verses 24--37}. The passage was selected beginning at verse 26, but thematic analysis revealed that verses 24--25 form an inseparable introduction, and verses 34--37 form the thematic continuation. The complete thematic unit is \textbf{8:24--37}.

\section{Thematic Structure}
The passage contains five interconnected themes:
\begin{enumerate}
\item \textbf{The Call to Life and Warning of Fitnah (8:24--25)}
\item \textbf{Remembrance of Divine Favour and Warning against Betrayal (8:26--28)}
\item \textbf{Taqwa and the Criterion (8:29)}
\item \textbf{The Quraysh's Plots, Arrogance, and Allah's Protection (8:30--33)}
\item \textbf{The Disbelievers' Deserved Punishment and Futile Spending (8:34--37)}
\end{enumerate}

\section{Sources}
All data drawn from the \textbf{Quran Wiki raw corpus} at \texttt{/root/tarbiyyah/quran-wiki/}:
\begin{itemize}
\item \textbf{Arabic:} Uthmani script, dual-sourced (Al Quran Cloud + Quran.com)
\item \textbf{Translations:} Sahih International, Pickthall, Yusuf Ali, Asad, Basmeih (Malay)
\item \textbf{Tafsir:} Ibn Kathir (Abridged) + Tafsir Muyassar
\end{itemize}

\chapter{Verse-by-Verse Analysis}

\section{Verse 8:24}
{\small\color{accent2}\itshape Theme 1: The Call to Life (8:24--25)}


\begin{quranbox}[Qur'an 8:24]
\quranverse{يَـٰٓأَيُّهَا ٱلَّذِينَ ءَامَنُوا۟ ٱسْتَجِيبُوا۟ لِلَّهِ وَلِلرَّسُولِ إِذَا دَعَاكُمْ لِمَا يُحْيِيكُمْ ۖ وَٱعْلَمُوٓا۟ أَنَّ ٱللَّهَ يَحُولُ بَيْنَ ٱلْمَرْءِ وَقَلْبِهِۦ وَأَنَّهُۥٓ إِلَيْهِ تُحْشَرُونَ}

\vspace{6pt}
\noindent\textbf{Sahih International:} O you who have believed, respond to Allah and to the Messenger when he calls you to that which gives you life. And know that Allah intervenes between a man and his heart and that to Him you will be gathered.

\vspace{4pt}
\noindent\textbf{Pickthall:} O ye who believe! Obey Allah, and the messenger when He calleth you to that which quickeneth you, and know that Allah cometh in between the man and his own heart, and that He it is unto Whom ye will be gathered.

\vspace{4pt}
\noindent\textbf{Yusuf Ali:} O ye who believe! give your response to Allah and His Messenger, when He calleth you to that which will give you life; and know that Allah cometh in between a man and his heart, and that it is He to Whom ye shall (all) be gathered.

\vspace{4pt}
\noindent\textbf{Asad:} O you who have attained to faith! Respond to the call of God and the Apostle whenever he calls you unto that which will give you life; and know that God intervenes between man and [the desires of] his heart, and that unto Him you shall be gathered.

\vspace{4pt}
\noindent\textbf{Basmeih (Malay):} Wahai orang-orang yang beriman, sahut dan sambutlah seruan Allah dan seruan RasulNya apabila Ia menyeru kamu kepada perkara-perkara yang menjadikan kamu hidup sempurna; dan ketahuilah bahawa sesungguhnya Allah berkuasa mengubah atau menyekat di antara seseorang itu dengan (pekerjaan) hatinya, dan sesungguhnya kepadaNyalah kamu akan dihimpunkan.
\end{quranbox}


\begin{tafsirbox}[Tafsir Ibn Kathir (Abridged)]
The Command to answer and obey Allah and His Messenger Al-Bukhari said, \textAR{اسْتَجِيبُواْ} "(Answer), obey, \textAR{لِمَا} \textAR{يُحْيِيكُمْ} (that which will give you life) that which will make your affairs good." Al-Bukhari went on to narrate that Abu Sa`id bin Al-Mu`alla said, "I was praying when the Prophet passed by and called me, but I did not answer him until I finished the prayer. He said, \textAR{«مَا} \textAR{مَنَعَكَ} \textAR{أَنْ} \textAR{تَأْتِيَنِي؟} \textAR{أَلَمْ} \textAR{يَقُلِ} \textAR{اللَّهُ}: (What prevented you from answering me Has not Allah said: \textAR{يأَيُّهَا} \textAR{الَّذِينَ} \textAR{ءَامَنُواْ} \textAR{اسْتَجِيبُواْ} \textAR{لِلَّهِ} \textAR{وَلِلرَّسُولِ} \textAR{إِذَا} \textAR{دَعَاكُمْ} \textAR{لِمَا} \textAR{يُحْيِيكُمْ} (O you who believe! Answer Allah and (His) Messenger when he calls you to that which will give you life)' He then said: \textAR{«لَأُعَلِّمَنَّكَ} \textAR{أَعْظَمَ} \textAR{سُورَةٍ} \textAR{فِي} \textAR{الْقُرآنِ} \textAR{قَبْلَ} \textAR{أَنْ} \textAR{أَخْرُج»} (I will teach you the greatest Suah in the Qur'an before I leave.) When he was about to leave, I mentioned what he said to me. He said, \textAR{الْحَمْدُ} \textAR{للَّهِ} \textAR{رَبِّ} \textAR{الْعَـلَمِينَ} (All the praises and thanks are to Allah, the Lord of all that exists...) 1:1-7. \textAR{«هِيَ} \textAR{السَّبْعُ} \textAR{الْمَثَانِي»} (Surely, it is the seven oft-repeated verses.)"' Muhammad bin Ishaq narrated that Muhammad bin Ja`far bin Az-Zubayr said that `Urwah bin Az-Zubayr explained this Ayah, \textAR{يأَيُّهَا} \textAR{الَّذِينَ} \textAR{ءَامَنُواْ} \textAR{اسْتَجِيبُواْ} \textAR{لِلَّهِ} \textAR{وَلِلرَّسُولِ} \textAR{إِذَا} \textAR{دَعَاكُمْ} \textAR{لِمَا} \textAR{يُحْيِيكُمْ} (O you who believe! Answer Allah and (His) Messenger when he calls you to that which will give you life,) "Answer when called to war (Jihad) with which Allah gives you might after meekness, and strength after weakness, and shields you from the enemy who oppressed you." Allah comes in between a Person and His Heart Allah said, \textAR{وَاعْلَمُواْ} \textAR{أَنَّ} \textAR{اللَّهَ} \textAR{يَحُولُ} \textAR{بَيْنَ} \textAR{الْمَرْءِ} \textAR{وَقَلْبِهِ} (and know that Allah comes in between a person and his heart.) Ibn `Abbas commented, "Allah prevents the believer from disbelief and the disbeliever from faith." Al-Hakim recorded this in his Mustadrak and said, "It is Sahih and they did not record it." . Similar was said by Mujahid, Sa`id, `Ikrimah, Ad-Dahhak, Abu Salih `Atiyyah, Muqatil bin Hayyan and As-Suddi. In another report from Mujahid, he commented; \textAR{يَحُولُ} \textAR{بَيْنَ} \textAR{الْمَرْءِ} \textAR{وَقَلْبِهِ} (...comes in between a person and his heart.) "Leaves him without comprehension," As-Suddi said, "Prevents one self from his own heart, so he will neither believe nor disbelieve except by His leave." There are several Hadiths that conform with the meaning of this Ayah. For instance, Imam Ahmad recorded that Anas bin Malik said, "The Prophet used to often say these words, \textAR{«يَا} \textAR{مُقَلِّبَ} \textAR{الْقُلُوبِ} \textAR{ثَبِّتْ} \textAR{قَلْبِي} \textAR{عَلَى} \textAR{دِينِك»} (O You Who changes the hearts, make my heart firm on Your religion.) We said, `O Allah's Messenger! We believed in you and in what you brought us. Are you afraid for us' He said, \textAR{«نَعَمْ،} \textAR{إِنَّ} \textAR{الْقُلُوبَ} \textAR{بَيْنَ} \textAR{إِصْبَعَيْنِ} \textAR{مِنْ} \textAR{أَصَابِعِ} \textAR{اللهِ} \textAR{تَعَالَى} \textAR{يُقَلِّبُهَا»} (Yes, for the hearts are between two of Allah's Fingers, He changes them (as He wills).)" This is the same narration recorded by At-Tirmidhi in the Book of Qadar in his Jami' Sunan, and he said, "Hasan." Imam Ahmad recorded that An-Nawwas bin Sam`an Al-Kilabi said that he heard the Prophet saying, \textAR{«مَا} \textAR{مِنْ} \textAR{قَلْبٍ} \textAR{إِلَّا} \textAR{وَهُوَ} \textAR{بَيْنَ} \textAR{أُصْبُعَيْنِ} \textAR{مِنْ} \textAR{أَصَابِعِ} \textAR{الرَّحْمَنِ} \textAR{رَبِّ} \textAR{الْعَالَمِينَ} \textAR{إِذَا} \textAR{شَاءَ} \textAR{أَنْ} \textAR{يُقِيمَهُ} \textAR{أَقَامَهُ} \textAR{وَإِذَا} \textAR{شَاءَ} \textAR{أَنْ} \textAR{يُزِيغَهُ} \textAR{أَزَاغَه»} (Every heart is between two of the Fingers of the Most Beneficent (Allah), Lord of all that exists, if He wills, He makes it straight, and if He wills, He makes it stray.) And he said: \textAR{«يَا} \textAR{مُقَلِّبَ} \textAR{الْقُلُوبِ} \textAR{ثَبِّتْ} \textAR{قَلْبِي} \textAR{عَلَى} \textAR{دِينِك»} (O You Who changes the hearts! keep my heart firm on Your religion) And he would say; \textAR{«وَالْمِيزَانُ} \textAR{بِيَدِ} \textAR{الرَّحْمنِ} \textAR{يَخْفِضُهُ} \textAR{وَيَرْفَعُه»} (The Balance is in the Hand of Ar-Rahman, He raises and lowers it.) This was also recorded by An-Nasai and Ibn Majah.
\end{tafsirbox}


\begin{muyassarbox}[Tafsir Muyassar]
\textAR{يا أيها الذين صدِّقوا بالله ربًا وبمحمد نبيًا ورسولا استجيبوا لله وللرسول بالطاعة إذا دعاكم لما يحييكم من الحق، ففي الاستجابة إصلاح حياتكم في الدنيا والآخرة، واعلموا -أيها المؤمنون- أن الله تعالى هو المتصرف في جميع الأشياء، والقادر على أن يحول بين الإنسان وما يشتهيه قلبه، فهو سبحانه الذي ينبغي أن يستجاب له إذا دعاكم; إذ بيده ملكوت كل شيء، واعلموا أنكم تُجمعون ليوم لا ريب فيه، فيجازي كلا بما يستحق.}
\end{muyassarbox}


\newpage

\section{Verse 8:25}
{\small\color{accent2}\itshape Theme 1: The Call to Life (8:24--25)}


\begin{quranbox}[Qur'an 8:25]
\quranverse{وَٱتَّقُوا۟ فِتْنَةً لَّا تُصِيبَنَّ ٱلَّذِينَ ظَلَمُوا۟ مِنكُمْ خَآصَّةً ۖ وَٱعْلَمُوٓا۟ أَنَّ ٱللَّهَ شَدِيدُ ٱلْعِقَابِ}

\vspace{6pt}
\noindent\textbf{Sahih International:} And fear a trial which will not strike those who have wronged among you exclusively, and know that Allah is severe in penalty.

\vspace{4pt}
\noindent\textbf{Pickthall:} And guard yourselves against a chastisement which cannot fall exclusively on those of you who are wrong-doers, and know that Allah is severe in punishment.

\vspace{4pt}
\noindent\textbf{Yusuf Ali:} And fear tumult or oppression, which affecteth not in particular (only) those of you who do wrong: and know that Allah is strict in punishment.

\vspace{4pt}
\noindent\textbf{Asad:} And beware of that temptation to evil which does not befall only those among you who are bent on denying the truth, to the exclusion of others; and know that God is severe in retribution.

\vspace{4pt}
\noindent\textbf{Basmeih (Malay):} Dan jagalah diri kamu daripada (berlakunya) dosa (yang membawa bala bencana) yang bukan sahaja akan menimpa orang-orang yang zalim di antara kamu secara khusus (tetapi akan menimpa kamu secara umum). Dan ketahuilah bahawa Allah Maha berat azab seksaNya.
\end{quranbox}


\begin{tafsirbox}[Tafsir Ibn Kathir (Abridged)]
Warning against an encompassing Fitnah Allah warns His believing servants of a Fitnah, trial and test, that encompasses the wicked and those around them. Therefore, such Fitnah will not be restricted to the sinners and evildoers. Rather, it will reach the others if the sins are not stopped and prevented. Imam Ahmad recorded that Mutarrif said, "We asked Az-Zubayr, `O Abu `Abdullah! What brought you here (for the battle of Al-Jamal) You abandoned the Khalifah who was assassinated (`Uthman, may Allah be pleased with him) and then came asking for revenge for his blood' He said, `We recited at the time of the Messenger of Allah (saw), and Abu Bakr, `Umar and `Uthman, \textAR{وَاتَّقُواْ} \textAR{فِتْنَةً} \textAR{لاَّ} \textAR{تُصِيبَنَّ} \textAR{الَّذِينَ} \textAR{ظَلَمُواْ} \textAR{مِنكُمْ} \textAR{خَآصَّةً} (And fear the Fitnah (affliction and trial) which affects not in particular (only) those of you who do wrong,) We did not think that this Ayah was about us too, until it reached us as it did."' `Ali bin Abi Talhah reported that Ibn `Abbas said that the Ayah, \textAR{وَاتَّقُواْ} \textAR{فِتْنَةً} \textAR{لاَّ} \textAR{تُصِيبَنَّ} \textAR{الَّذِينَ} \textAR{ظَلَمُواْ} \textAR{مِنكُمْ} \textAR{خَآصَّةً} (And fear the Fitnah (affliction and trial) which affects not in particular (only) those of you who do wrong,) refers to the Companions of the Prophet in particular. In another narration from Ibn `Abbas, he said, "Allah commanded the believers to stop evil from flourishing among them, so that Allah does not encompass them all in the torment (Fitnah). " This, indeed, is a very good explanation, prompting Mujahid to comment about Allah's statement, \textAR{وَاتَّقُواْ} \textAR{فِتْنَةً} \textAR{لاَّ} \textAR{تُصِيبَنَّ} \textAR{الَّذِينَ} \textAR{ظَلَمُواْ} \textAR{مِنكُمْ} \textAR{خَآصَّةً} (And fear the Fitnah (affliction and trial) which affects not in particular (only) those of you who do wrong,) "Is for you too!" Several said similarly, such as Ad-Dahhak and Yazid... [truncated]
\end{tafsirbox}


\begin{muyassarbox}[Tafsir Muyassar]
\textAR{واحذروا -أيها المؤمنون- اختبارًا ومحنة يُعَمُّ بها المسيء وغيره لا يُخَص بها أهل المعاصي ولا مَن باشر الذنب، بل تصيب الصالحين معهم إذا قدروا على إنكار الظلم ولم ينكروه، واعلموا أن الله شديد العقاب لمن خالف أمره ونهيه.}
\end{muyassarbox}


\newpage

\section{Verse 8:26}
{\small\color{accent2}\itshape Theme 2: Divine Favour and Betrayal (8:26--28)}


\begin{quranbox}[Qur'an 8:26]
\quranverse{وَٱذْكُرُوٓا۟ إِذْ أَنتُمْ قَلِيلٌ مُّسْتَضْعَفُونَ فِى ٱلْأَرْضِ تَخَافُونَ أَن يَتَخَطَّفَكُمُ ٱلنَّاسُ فَـَٔاوَىٰكُمْ وَأَيَّدَكُم بِنَصْرِهِۦ وَرَزَقَكُم مِّنَ ٱلطَّيِّبَـٰتِ لَعَلَّكُمْ تَشْكُرُونَ}

\vspace{6pt}
\noindent\textbf{Sahih International:} And remember when you were few and oppressed in the land, fearing that people might abduct you, but He sheltered you, supported you with His victory, and provided you with good things - that you might be grateful.

\vspace{4pt}
\noindent\textbf{Pickthall:} And remember, when ye were few and reckoned feeble in the land, and were in fear lest men should extirpate you, how He gave you refuge, and strengthened you with His help, and made provision of good things for you, that haply ye might be thankful.

\vspace{4pt}
\noindent\textbf{Yusuf Ali:} Call to mind when ye were a small (band), despised through the land, and afraid that men might despoil and kidnap you; But He provided a safe asylum for you, strengthened you with His aid, and gave you Good things for sustenance: that ye might be grateful.

\vspace{4pt}
\noindent\textbf{Asad:} And remember the time when you were few [and] helpless on earth, fearful lest people do away with you - whereupon He sheltered you, and strengthened you with His succour, and provided for you sustenance out of the good things of life, so that you might have cause to be grateful.

\vspace{4pt}
\noindent\textbf{Basmeih (Malay):} Dan ingatlah ketika kamu sedikit bilangannya serta tertindas di bumi, kamu takut orang-orang menangkap dan melarikan kamu, maka Allah memberi kamu tempat bermustautin dan diperkuatkanNya kamu dengan pertolonganNya, serta dikurniakanNya kamu dari rezeki yang baik-baik, supaya kamu bersyukur.
\end{quranbox}


\begin{tafsirbox}[Tafsir Ibn Kathir (Abridged)]
Reminding Muslims of Their previous State of Weakness and Subjugation which changed into Might and Triumph Allah, the Exalted, reminds His believing servants of His blessings and favors on them. They were few and He made them many, weak and fearful and He provided them with strength and victory. They were meek and poor, and He granted them sustenance and livelihood. He ordered them to be grateful to Him, and they obeyed Him and implemented what He commanded. When the believers were still in Makkah they were few, practicing their religion in secret, oppressed, fearing that pagans, fire worshippers or Romans might kidnap them from the various parts of Allah's earth, for they were all enemies of the Muslims, especially since Muslims were few and weak. Later on, Allah permitted the believers to migrate to Al-Madinah, where He allowed them to settle in a safe resort. Allah made the people of Al-Madinah their allies, giving them refuge and support during Badr and other battles. They helped the Migrants with their wealth and gave up their lives in obedience of Allah and His Messenger . Qatadah bin Di`amah As-Sadusi commented, \textAR{وَاذْكُرُواْ} \textAR{إِذْ} \textAR{أَنتُمْ} \textAR{قَلِيلٌ} \textAR{مُّسْتَضْعَفُونَ} \textAR{فِى} \textAR{الاٌّرْضِ} (And remember when you were few and were reckoned weak in the land,) "Arabs were the weakest of the weak, had the toughest life, the emptiest stomachs, the barest skin and the most obvious misguidance. Those who lived among them lived in misery; those who died went to the Fire. They were being eaten up, but unable to eat up others! By Allah! We do not know of a people on the face of the earth at that time who had a worse life than them. When Allah brought Islam, He made it dominant on the earth, thus bringing provisions and leadership for them over the necks of people. It is through Islam that Allah granted all what you see, so thank Him for His favors, for your Lord is One Who bestows favors and likes praise. Verily, those w... [truncated]
\end{tafsirbox}


\begin{muyassarbox}[Tafsir Muyassar]
\textAR{واذكروا أيها المؤمنون نِعَم الله عليكم إذ أنتم بـ «مكة» قليلو العدد مقهورون، تخافون أن يأخذكم الكفار بسرعة، فجعل لكم مأوى تأوون إليه وهو «المدينة» ، وقوَّاكم بنصره عليهم يوم «بدر» ، وأطعمكم من الطيبات -التي من جملتها الغنائم-; لكي تشكروا له على ما رزقكم وأنعم به عليكم.}
\end{muyassarbox}


\newpage

\section{Verse 8:27}
{\small\color{accent2}\itshape Theme 2: Divine Favour and Betrayal (8:26--28)}


\begin{quranbox}[Qur'an 8:27]
\quranverse{يَـٰٓأَيُّهَا ٱلَّذِينَ ءَامَنُوا۟ لَا تَخُونُوا۟ ٱللَّهَ وَٱلرَّسُولَ وَتَخُونُوٓا۟ أَمَـٰنَـٰتِكُمْ وَأَنتُمْ تَعْلَمُونَ}

\vspace{6pt}
\noindent\textbf{Sahih International:} O you who have believed, do not betray Allah and the Messenger or betray your trusts while you know [the consequence].

\vspace{4pt}
\noindent\textbf{Pickthall:} O ye who believe! Betray not Allah and His messenger, nor knowingly betray your trusts.

\vspace{4pt}
\noindent\textbf{Yusuf Ali:} O ye that believe! betray not the trust of Allah and the Messenger, nor misappropriate knowingly things entrusted to you.

\vspace{4pt}
\noindent\textbf{Asad:} [Hence,] O you who have attained to faith, do not be false to God and the Apostle, and do not knowingly be false to the trust that has been reposed in you;

\vspace{4pt}
\noindent\textbf{Basmeih (Malay):} Wahai orang-orang yang beriman! Janganlah kamu mengkhianati (amanah) Allah dan RasulNya, dan (janganlah) kamu mengkhianati amanah-amanah kamu, sedang kamu mengetahui (salahnya).
\end{quranbox}


\begin{tafsirbox}[Tafsir Ibn Kathir (Abridged)]
Reason behind revealing This Ayah, and the prohibition of Betrayal The Two Sahihs mention the story of Hatib bin Abi Balta`ah. In the year of the victory of Makkah he wrote to the Quraysh alerting them that the Messenger of Allah (saw) intended to march towards them. Allah informed His Messenger of this, and he sent a Companion to retrieve the letter that Hatib sent, and then he summoned him. He admitted to what he did. `Umar bin Al-Khattab stood up and said, "O Allah's Messenger! Should I cut off his head, for he has betrayed Allah, His Messenger and the believers" The Prophet said, \textAR{«دَعْهُ} \textAR{فَإِنَّهُ} \textAR{قَدْ} \textAR{شَهِدَ} \textAR{بَدْرًا،} \textAR{وَمَا} \textAR{يُدْرِيكَ} \textAR{لَعَلَّ} \textAR{اللهَ} \textAR{اطَّلَعَ} \textAR{عَلَى} \textAR{أَهْلِ} \textAR{بَدْرٍ} \textAR{فَقَالَ}: \textAR{اعْمَلُوا} \textAR{مَا} \textAR{شِئْتُمْ} \textAR{فَقَدْ} \textAR{غَفَرْتُ} \textAR{لَكُم»} (Leave him! He participated in Badr. How do you know that Allah has not looked at those who participated in Badr and said, Do whatever you want, for I have forgiven you.) However, it appears that this Ayah is more general, even if it was revealed about a specific incident. Such rulings are dealt with by their indications, not the specific reasons behind revealing them, according to the majority of scholars. Betrayal includes both minor and major sins, as well those that affect others. `Ali bin Abi Talhah said that Ibn `Abbas commented on the Ayah, \textAR{وَتَخُونُواْ} \textAR{أَمَـنَـتِكُمْ} (nor betray your Amanat) "The Amanah refers to the actions that Allah has entrusted the servants with, such as and including what He ordained. Therefore, Allah says here, \textAR{لاَ} \textAR{تَخُونُواْ} (nor betray...), `do not abandon the obligations."' `Abdur-Rahman bin Zayd commented, "Allah forbade you from betraying Him and His Messenger, as hypocrites do." Allah said, \textAR{وَاعْلَمُواْ} \textAR{أَنَّمَآ} \textAR{أَمْوَلُكُمْ} \textAR{وَأَوْلَـدُكُمْ} \textAR{فِتْنَةٌ} (And know that your possessions and your children are but a trial.) from Him to you. He grants these to you so that He knows which of you will be grateful and obedient to Him, or become busy with and dedicated to them instead of Him. Allah said in another Ayah, \textAR{إِنَّمَآ} \textAR{أَمْوَلُكُمْ} \textAR{وَأَوْلَـدُكُمْ} \textAR{فِتْنَةٌ} \textAR{وَاللَّهُ} \textAR{عِنْدَهُ} \textAR{أَجْرٌ} \textAR{عَظِيمٌ} (Your wealth and your children are only a trial, whereas Allah! With Him is a great reward.) 64:15, \textAR{وَنَبْلُوكُم} \textAR{بِالشَّرِّ} \textAR{وَالْخَيْرِ} \textAR{فِتْنَةً} (And We shall make a trial of you with evil and with good.) 21:35, \textAR{يأَيُّهَا} \textAR{الَّذِينَ} \textAR{ءَامَنُواْ} \textAR{لاَ} \textAR{تُلْهِكُمْ} \textAR{أَمْوَلُكُمْ} \textAR{وَلاَ} \textAR{أَوْلَـدُكُمْ} \textAR{عَن} \textAR{ذِكْرِ} \textAR{اللَّهِ} \textAR{وَمَن} \textAR{يَفْعَلْ} \textAR{ذَلِكَ} \textAR{فَأُوْلَـئِكَ} \textAR{هُمُ} \textAR{الْخَـسِرُونَ} (O you who believe! Let not your properties or your children divert you from the remembrance of Allah. And whosoever does that, then they are the losers.)63:9, and, \textAR{يأَيُّهَا} \textAR{الَّذِينَ} \textAR{ءَامَنُواْ} \textAR{إِنَّ} \textAR{مِنْ} \textAR{أَزْوَجِكُمْ} \textAR{وَأَوْلـدِكُمْ} \textAR{عَدُوّاً} \textAR{لَّكُمْ} \textAR{فَاحْذَرُوهُمْ} (O you who believe! Verily, among your wives and your children there are enemies for you (who may stop you from the obedience of Allah); therefore beware of them!) 64:14 Allah said next, \textAR{وَأَنَّ} \textAR{اللَّهَ} \textAR{عِندَهُ} \textAR{أَجْرٌ} \textAR{عَظِيمٌ} (And that surely with Allah is a mighty reward.) Therefore, Allah's reward, favor and Paradise are better for you than wealth and children. Certainly, among the wealth and children there might be enemies for you and much of them avail nothing. With Allah alone is the decision and sovereignty in this life and the Hereafter, and He gives tremendous rewards on the Day of Resurrection. In the Sahih, there is a Hadith in which the Messenger of Allah (saw) said, \textAR{«ثَلَاثٌ} \textAR{مَنْ} \textAR{كُنَّ} \textAR{فِيهِ،} \textAR{وَجَدَ} \textAR{بِهِنَّ} \textAR{حَلَاوَةَ} \textAR{الْإِيمَانِ}: \textAR{مَنْ} \textAR{كَانَ} \textAR{اللهُ} \textAR{وَرَسُولُهُ} \textAR{أَحَبَّ} \textAR{إِلَيْهِ} \textAR{مِمَّا} \textAR{سِوَاهُمَا،} \textAR{وَمَنْ} \textAR{كَانَ} \textAR{يُحِبُّ} \textAR{الْمَرْءَ} \textAR{لَا} \textAR{يُحِبُّهُ} \textAR{إِلَّا} \textAR{للهِ،} \textAR{وَمَنْ} \textAR{كَانَ} \textAR{أَنْ} \textAR{يُلْقَى} \textAR{فِي} \textAR{النَّارِ} \textAR{أَحَبَّ} \textAR{إِلَيْهِ} \textAR{مِنْ} \textAR{أَنْ} \textAR{يَرْجِعَ} \textAR{إِلَى} \textAR{الْكُفْرِ} \textAR{بَعْدَ} \textAR{إِذْ} \textAR{أَنْقَذَهُ} \textAR{اللهُ} \textAR{مِنْه»} (There are three qualities for which whomever has them, he will have tasted the sweetness of faith. (They are:) whoever Allah and His Messenger are dearer to him than anyone els [...] [Tafsir continues in full corpus]
\end{tafsirbox}


\begin{muyassarbox}[Tafsir Muyassar]
\textAR{يا أيها الذين صدَّقوا الله ورسوله وعملوا بشرعه لا تخونوا الله ورسوله بترك ما أوجبه الله عليكم وفِعْل ما نهاكم عنه، ولا تفرطوا فيما ائتمنكم الله عليه، وأنتم تعلمون أنه أمانة يجب الوفاء بها.}
\end{muyassarbox}


\newpage

\section{Verse 8:28}
{\small\color{accent2}\itshape Theme 2: Divine Favour and Betrayal (8:26--28)}


\begin{quranbox}[Qur'an 8:28]
\quranverse{وَٱعْلَمُوٓا۟ أَنَّمَآ أَمْوَٰلُكُمْ وَأَوْلَـٰدُكُمْ فِتْنَةٌ وَأَنَّ ٱللَّهَ عِندَهُۥٓ أَجْرٌ عَظِيمٌ}

\vspace{6pt}
\noindent\textbf{Sahih International:} And know that your properties and your children are but a trial and that Allah has with Him a great reward.

\vspace{4pt}
\noindent\textbf{Pickthall:} And know that your possessions and your children are a test, and that with Allah is immense reward.

\vspace{4pt}
\noindent\textbf{Yusuf Ali:} And know ye that your possessions and your progeny are but a trial; and that it is Allah with Whom lies your highest reward.

\vspace{4pt}
\noindent\textbf{Asad:} and know that your worldly goods and your children are but a trial and a temptation, and that with God there is a tremendous reward.

\vspace{4pt}
\noindent\textbf{Basmeih (Malay):} Dan ketahuilah bahawa harta benda kamu dan anak-anak kamu itu hanyalah menjadi ujian, dan sesungguhnya di sisi Allah jualah pahala yang besar.
\end{quranbox}


\begin{tafsirbox}[Tafsir Ibn Kathir (Abridged)]
Reason behind revealing This Ayah, and the prohibition of Betrayal The Two Sahihs mention the story of Hatib bin Abi Balta`ah. In the year of the victory of Makkah he wrote to the Quraysh alerting them that the Messenger of Allah (saw) intended to march towards them. Allah informed His Messenger of this, and he sent a Companion to retrieve the letter that Hatib sent, and then he summoned him. He admitted to what he did. `Umar bin Al-Khattab stood up and said, "O Allah's Messenger! Should I cut off his head, for he has betrayed Allah, His Messenger and the believers" The Prophet said, \textAR{«دَعْهُ} \textAR{فَإِنَّهُ} \textAR{قَدْ} \textAR{شَهِدَ} \textAR{بَدْرًا،} \textAR{وَمَا} \textAR{يُدْرِيكَ} \textAR{لَعَلَّ} \textAR{اللهَ} \textAR{اطَّلَعَ} \textAR{عَلَى} \textAR{أَهْلِ} \textAR{بَدْرٍ} \textAR{فَقَالَ}: \textAR{اعْمَلُوا} \textAR{مَا} \textAR{شِئْتُمْ} \textAR{فَقَدْ} \textAR{غَفَرْتُ} \textAR{لَكُم»} (Leave him! He participated in Badr. How do you know that Allah has not looked at those who participated in Badr and said, Do whatever you want, for I have forgiven you.) However, it appears that this Ayah is more general, even if it was revealed about a specific incident. Such rulings are dealt with by their indications, not the specific reasons behind revealing them, according to the majority of scholars. Betrayal includes both minor and major sins, as well those that affect others. `Ali bin Abi Talhah said that Ibn `Abbas commented on the Ayah, \textAR{وَتَخُونُواْ} \textAR{أَمَـنَـتِكُمْ} (nor betray your Amanat) "The Amanah refers to the actions that Allah has entrusted the servants with, such as and including what He ordained. Therefore, Allah says here, \textAR{لاَ} \textAR{تَخُونُواْ} (nor betray...), `do not abandon the obligations."' `Abdur-Rahman bin Zayd commented, "Allah forbade you from betraying Him and His Messenger, as hypocrites do." Allah said, \textAR{وَاعْلَمُواْ} \textAR{أَنَّمَآ} \textAR{أَمْوَلُكُمْ} \textAR{وَأَوْلَـدُكُمْ} \textAR{فِتْنَةٌ} (And know that your possessions and your children are but a trial.) from Him to you. He grants these to you so that He knows which of you will be grateful and obedient to Him, or become busy with and dedicated to them instead of Him. Allah said in another Ayah, \textAR{إِنَّمَآ} \textAR{أَمْوَلُكُمْ} \textAR{وَأَوْلَـدُكُمْ} \textAR{فِتْنَةٌ} \textAR{وَاللَّهُ} \textAR{عِنْدَهُ} \textAR{أَجْرٌ} \textAR{عَظِيمٌ} (Your wealth and your children are only a trial, whereas Allah! With Him is a great reward.) 64:15, \textAR{وَنَبْلُوكُم} \textAR{بِالشَّرِّ} \textAR{وَالْخَيْرِ} \textAR{فِتْنَةً} (And We shall make a trial of you with evil and with good.) 21:35, \textAR{يأَيُّهَا} \textAR{الَّذِينَ} \textAR{ءَامَنُواْ} \textAR{لاَ} \textAR{تُلْهِكُمْ} \textAR{أَمْوَلُكُمْ} \textAR{وَلاَ} \textAR{أَوْلَـدُكُمْ} \textAR{عَن} \textAR{ذِكْرِ} \textAR{اللَّهِ} \textAR{وَمَن} \textAR{يَفْعَلْ} \textAR{ذَلِكَ} \textAR{فَأُوْلَـئِكَ} \textAR{هُمُ} \textAR{الْخَـسِرُونَ} (O you who believe! Let not your properties or your children divert you from the remembrance of Allah. And whosoever does that, then they are the losers.)63:9, and, \textAR{يأَيُّهَا} \textAR{الَّذِينَ} \textAR{ءَامَنُواْ} \textAR{إِنَّ} \textAR{مِنْ} \textAR{أَزْوَجِكُمْ} \textAR{وَأَوْلـدِكُمْ} \textAR{عَدُوّاً} \textAR{لَّكُمْ} \textAR{فَاحْذَرُوهُمْ} (O you who believe! Verily, among your wives and your children there are enemies for you (who may stop you from the obedience of Allah); therefore beware of them!) 64:14 Allah said next, \textAR{وَأَنَّ} \textAR{اللَّهَ} \textAR{عِندَهُ} \textAR{أَجْرٌ} \textAR{عَظِيمٌ} (And that surely with Allah is a mighty reward.) Therefore, Allah's reward, favor and Paradise are better for you than wealth and children. Certainly, among the wealth and children there might be enemies for you and much of them avail nothing. With Allah alone is the decision and sovereignty in this life and the Hereafter, and He gives tremendous rewards on the Day of Resurrection. In the Sahih, there is a Hadith in which the Messenger of Allah (saw) said, \textAR{«ثَلَاثٌ} \textAR{مَنْ} \textAR{كُنَّ} \textAR{فِيهِ،} \textAR{وَجَدَ} \textAR{بِهِنَّ} \textAR{حَلَاوَةَ} \textAR{الْإِيمَانِ}: \textAR{مَنْ} \textAR{كَانَ} \textAR{اللهُ} \textAR{وَرَسُولُهُ} \textAR{أَحَبَّ} \textAR{إِلَيْهِ} \textAR{مِمَّا} \textAR{سِوَاهُمَا،} \textAR{وَمَنْ} \textAR{كَانَ} \textAR{يُحِبُّ} \textAR{الْمَرْءَ} \textAR{لَا} \textAR{يُحِبُّهُ} \textAR{إِلَّا} \textAR{للهِ،} \textAR{وَمَنْ} \textAR{كَانَ} \textAR{أَنْ} \textAR{يُلْقَى} \textAR{فِي} \textAR{النَّارِ} \textAR{أَحَبَّ} \textAR{إِلَيْهِ} \textAR{مِنْ} \textAR{أَنْ} \textAR{يَرْجِعَ} \textAR{إِلَى} \textAR{الْكُفْرِ} \textAR{بَعْدَ} \textAR{إِذْ} \textAR{أَنْقَذَهُ} \textAR{اللهُ} \textAR{مِنْه»} (There are three qualities for which whomever has them, he will have tasted the sweetness of faith. (They are:) whoever Allah and His Messenger are dearer to him than anyone els [...] [Tafsir continues in full corpus]
\end{tafsirbox}


\begin{muyassarbox}[Tafsir Muyassar]
\textAR{واعلموا -أيها المؤمنون- أن أموالكم التي استخلفكم الله فيها، وأولادكم الذين وهبهم الله لكم اختبار من الله وابتلاء لعباده; ليعلم أيشكرونه عليها ويطيعونه فيها، أو ينشغلون بها عنه؟ واعلموا أن الله عنده خير وثواب عظيم لمن اتقاه وأطاعه.}
\end{muyassarbox}


\newpage

\section{Verse 8:29}
{\small\color{accent2}\itshape Theme 3: Taqwa and the Criterion (8:29)}


\begin{quranbox}[Qur'an 8:29]
\quranverse{يَـٰٓأَيُّهَا ٱلَّذِينَ ءَامَنُوٓا۟ إِن تَتَّقُوا۟ ٱللَّهَ يَجْعَل لَّكُمْ فُرْقَانًا وَيُكَفِّرْ عَنكُمْ سَيِّـَٔاتِكُمْ وَيَغْفِرْ لَكُمْ ۗ وَٱللَّهُ ذُو ٱلْفَضْلِ ٱلْعَظِيمِ}

\vspace{6pt}
\noindent\textbf{Sahih International:} O you who have believed, if you fear Allah, He will grant you a criterion and will remove from you your misdeeds and forgive you. And Allah is the possessor of great bounty.

\vspace{4pt}
\noindent\textbf{Pickthall:} O ye who believe! If ye keep your duty to Allah, He will give you discrimination (between right and wrong) and will rid you of your evil thoughts and deeds, and will forgive you. Allah is of Infinite Bounty.

\vspace{4pt}
\noindent\textbf{Yusuf Ali:} O ye who believe! if ye fear Allah, He will grant you a criterion (to judge between right and wrong), remove from you (all) evil (that may afflict) you, and forgive you: for Allah is the Lord of grace unbounded.

\vspace{4pt}
\noindent\textbf{Asad:} O you who have attained to faith! If you remain conscious of God. He will endow you with a standard by which to discern the true from the false, and will efface your bad deeds, and will forgive you your sins: for God is limitless in His great bounty.

\vspace{4pt}
\noindent\textbf{Basmeih (Malay):} Wahai orang-orang yang beriman! Jika kamu bertaqwa kepada Allah, nescaya Ia mengadakan bagi kamu (petunjuk) yang membezakan antara yang benar dengan yang salah, dan menghapuskan kesalahan-kesalahan kamu, serta mengampunkan (dosa-dosa) kamu. Dan Allah (sememangnya) mempunyai limpah kurnia yang besar.
\end{quranbox}


\begin{tafsirbox}[Tafsir Ibn Kathir (Abridged)]
Ibn `Abbas, As-Suddi, Mujahid, `Ikrimah, Ad-Dahhak, Qatadah, Muqatil bin Hayyan and several others said that, \textAR{فُرْقَانًا} (Furqan), means, `a way out'; Mujahid added, "In this life and the Hereafter." In another narration, Ibn `Abbas is reported to have said, `Furqan' means `salvation' or -- according to another narration -- `aid'. Muhammad bin Ishaq said that `Furqan' means `criterion between truth and falsehood'. This last explanation from Ibn Ishaq is more general than the rest that we mentioned, and it also includes the other meanings. Certainly, those who have Taqwa of Allah by obeying what He ordained and abstaining from what he forbade, will be guided to differentiate between the truth and the falsehood. This will be a triumph, safety and a way out for them from the affairs of this life, all the while acquiring happiness in the Hereafter. They will also gain forgiveness, thus having their sins erased, and pardon, thus having their sins covered from other people, as well as, being directed to a way to gain Allah's tremendous rewards, \textAR{يأَيُّهَا} \textAR{الَّذِينَ} \textAR{ءَامَنُواْ} \textAR{اتَّقُواْ} \textAR{اللَّهَ} \textAR{وَءَامِنُواْ} \textAR{بِرَسُولِهِ} \textAR{يُؤْتِكُمْ} \textAR{كِفْلَيْنِ} \textAR{مِن} \textAR{رَّحْمَتِهِ} \textAR{وَيَجْعَل} \textAR{لَّكُمْ} \textAR{نُوراً} \textAR{تَمْشُونَ} \textAR{بِهِ} \textAR{وَيَغْفِرْ} \textAR{لَكُمْ} \textAR{وَاللَّهُ} \textAR{غَفُورٌ} \textAR{رَّحِيمٌ} (O you who believe! Have Taqwa of Allah, and believe in His Messenger, He will give you a double portion of His mercy, and He will give you a light by which you shall walk (straight). And He will forgive you. And Allah is Oft-Forgiving, Most Merciful.) 57:28.
\end{tafsirbox}


\begin{muyassarbox}[Tafsir Muyassar]
\textAR{يا أيها الذين صدَّقوا الله ورسوله وعملوا بشرعه إن تتقوا الله بفعل أوامره واجتناب نواهيه يجعل لكم فصلا بين الحق والباطل، ويَمحُ عنكم ما سلف من ذنوبكم ويسترها عليكم، فلا يؤاخذكم بها. والله ذو الإحسان والعطاء الكثير الواسع.}
\end{muyassarbox}


\newpage

\section{Verse 8:30}
{\small\color{accent2}\itshape Theme 4: Quraysh's Plots and Allah's Protection (8:30--33)}


\begin{quranbox}[Qur'an 8:30]
\quranverse{وَإِذْ يَمْكُرُ بِكَ ٱلَّذِينَ كَفَرُوا۟ لِيُثْبِتُوكَ أَوْ يَقْتُلُوكَ أَوْ يُخْرِجُوكَ ۚ وَيَمْكُرُونَ وَيَمْكُرُ ٱللَّهُ ۖ وَٱللَّهُ خَيْرُ ٱلْمَـٰكِرِينَ}

\vspace{6pt}
\noindent\textbf{Sahih International:} And [remember, O Muhammad], when those who disbelieved plotted against you to restrain you or kill you or evict you [from Makkah]. But they plan, and Allah plans. And Allah is the best of planners.

\vspace{4pt}
\noindent\textbf{Pickthall:} And when those who disbelieve plot against thee (O Muhammad) to wound thee fatally, or to kill thee or to drive thee forth; they plot, but Allah (also) plotteth; and Allah is the best of plotters.

\vspace{4pt}
\noindent\textbf{Yusuf Ali:} Remember how the Unbelievers plotted against thee, to keep thee in bonds, or slay thee, or get thee out (of thy home). They plot and plan, and Allah too plans; but the best of planners is Allah.

\vspace{4pt}
\noindent\textbf{Asad:} AND [remember, O Prophet,] how those who were bent on denying the truth were scheming against thee, in order to restrain thee [from preaching], or to slay thee, or to drive thee away: thus have they [always] schemed: but God brought their scheming to nought-for God is above all schemers.

\vspace{4pt}
\noindent\textbf{Basmeih (Malay):} Dan ingatlah (wahai Muhammad), ketika orang-orang kafir musyrik (Makkah) menjalankan tipu daya terhadapmu untuk menahanmu, atau membunuhmu, atau mengusirmu. Mereka menjalankan tipu daya dan Allah menggagalkan tipu daya (mereka), kerana Allah sebaik-baik yang menggagalkan tipu daya.
\end{quranbox}


\begin{tafsirbox}[Tafsir Ibn Kathir (Abridged)]
The Makkans plot to kill the Prophet , imprison Him or expel Him from Makkah Ibn `Abbas, Mujahid and Qatadah said, \textAR{لِيُثْبِتُوكَ} (Liyuthbituka) means "to imprison you." As-Suddi said, "Ithbat is to confine or to shackle." Imam Muhammad bin Ishaq bin Yasar, the author of Al-Maghazi, reported from `Abdullah bin Abi Najih, from Mujahid, from Ibn `Abbas, "Some of the chiefs of the various tribes of Quraysh gathered in Dar An-Nadwah (their conference area) and Iblis (Shaytan) met them in the shape of an eminent old man. When they saw him, they asked, `Who are you' He said, `An old man from Najd. I heard that you are having a meeting, and I wished to attend your meeting. You will benefit from my opinion and advice.' They said, `Agreed, come in.' He entered with them. Iblis said, `You have to think about this man (Muhammad)! By Allah, he will soon overwhelm you with his matter (religion).' One of them said, `Imprison him, restrained in chains, until he dies just like the poets before him all died, such as Zuhayr and An-Nabighah! Verily, he is a poet like they were.' The old man from Najd, the enemy of Allah, commented, `By Allah! This is not a good idea. His Lord will release him from his prison to his companions, who will liberate him from your hands. They will protect him from you and they might expel you from your land.' They said, `This old man said the truth. Therefore, seek an opinion other than this one.' Another one of them said, `Expel him from your land, so that you are free from his trouble! If he leaves your land, you will not be bothered by what he does or where he goes, as long as he is not among you to bring you troubles, he will be with someone else.' The old man from Najd replied, `By Allah! This is not a good opinion. Have you forgotten his sweet talk and eloquency, as well as, how his speech captures the hearts By Allah! This way, he will collect even more followers among Arabs, who will gather against you and attack you in your own land, expel you and kill your chiefs.' They said, `He has said the truth, by Allah! Therefore, seek an opinion other than this one.' hAbu Jahl, may Allah curse him, spoke next, `By Allah! I have an idea that no one else has suggested yet, and I see no better opinion for you. Choose a strong, socially elevated young man from each tribe, and give each one of them a sharp sword. Then they would all strike Muhammad at the same time with their swords and kill him. Hence, his blood would be shed by all tribes. This way, his tribe, Banu Hashim, would realize that they cannot wage war against all of the Quraysh tribes and would be forced to agree to accept the blood money; we would have brought comfort to ourselves and stopped him from bothering us.' The old man from Najd commented, `By Allah! This man has expressed the best opinion, and I do not support any other opinion.' They quickly ended their meeting and started preparing for the implementation of this plan. Jibril came to the Prophet and commanded him not to sleep in his bed that night and conveyed to him the news of their plot. The Messenger of Allah (saw) did not sleep in his house that night, and Allah gave him permission to migrate. After the Messenger (saw) migrated to Al-Madinah, Allah revealed to him Suat Al-Anfal reminding him of His favors and the bounties He gave him, \textAR{وَإِذْ} \textAR{يَمْكُرُ} \textAR{بِكَ} \textAR{الَّذِينَ} \textAR{كَفَرُواْ} \textAR{لِيُثْبِتُوكَ} \textAR{أَوْ} \textAR{يَقْتُلُوكَ} \textAR{أَوْ} \textAR{يُخْرِجُوكَ} \textAR{وَيَمْكُرُونَ} \textAR{وَيَمْكُرُ} \textAR{اللَّهُ} \textAR{وَاللَّهُ} \textAR{خَيْرُ} \textAR{الْمَـكِرِينَ} (And (remember) when the disbelievers plotted against you to imprison you, or to kill you, or to expel you (from Makkah); they were plotting and Allah too was plotting; and Allah is the best of plotters.) Allah replied to the pagans' statement that they should await the death of the Prophet , just as the poets before him perished, as they claimed, \textAR{أَمْ} \textAR{يَقُولُونَ} \textAR{شَاعِرٌ} \textAR{نَّتَرَبَّصُ} \textAR{بِهِ} \textAR{رَيْبَ} \textAR{الْمَنُونِ} (Or do th [...] [Tafsir continues in full corpus]
\end{tafsirbox}


\begin{muyassarbox}[Tafsir Muyassar]
\textAR{واذكر -أيها الرسول- حين يكيد لك مشركو قومك بـ «مكَّة» ; ليحبسوك أو يقتلوك أو ينفوك من بلدك. ويكيدون لك، وردَّ الله مكرهم عليهم جزاء لهم، ويمكر الله، والله خير الماكرين.}
\end{muyassarbox}


\newpage

\section{Verse 8:31}
{\small\color{accent2}\itshape Theme 4: Quraysh's Plots and Allah's Protection (8:30--33)}


\begin{quranbox}[Qur'an 8:31]
\quranverse{وَإِذَا تُتْلَىٰ عَلَيْهِمْ ءَايَـٰتُنَا قَالُوا۟ قَدْ سَمِعْنَا لَوْ نَشَآءُ لَقُلْنَا مِثْلَ هَـٰذَآ ۙ إِنْ هَـٰذَآ إِلَّآ أَسَـٰطِيرُ ٱلْأَوَّلِينَ}

\vspace{6pt}
\noindent\textbf{Sahih International:} And when Our verses are recited to them, they say, "We have heard. If we willed, we could say [something] like this. This is not but legends of the former peoples."

\vspace{4pt}
\noindent\textbf{Pickthall:} And when Our revelations are recited unto them they say: We have heard. If we wish we can speak the like of this. Lo! this is naught but fables of the men of old.

\vspace{4pt}
\noindent\textbf{Yusuf Ali:} When Our Signs are rehearsed to them, they say: "We have heard this (before): if we wished, we could say (words) like these: these are nothing but tales of the ancients."

\vspace{4pt}
\noindent\textbf{Asad:} And whenever Our messages were conveyed to them, they would say, "We have heard [all this] before; if we wanted, we could certainly compose sayings like these [ourselves]: they are nothing but fables of ancient times!"

\vspace{4pt}
\noindent\textbf{Basmeih (Malay):} Dan apabila dibacakan kepada mereka ayat-ayat Kami, mereka berkata: "Sesungguhnya kami telah mendengarnya. Kalau kami mahu, nescaya kami dapat mengatakan (kata-kata) seperti (Al-Quran) ini. (Al-Quran) ini tidak lain hanyalah cerita cerita dongeng orang-orang dahulu kala".
\end{quranbox}


\begin{tafsirbox}[Tafsir Ibn Kathir (Abridged)]
The Quraysh claimed They can produce Something similar to the Qur'an Allah describes the disbelief, transgression, rebellion, as well as misguided statements that the pagans of Quraysh used to utter when they heard Allah's Ayat being recited to them, \textAR{قَدْ} \textAR{سَمِعْنَا} \textAR{لَوْ} \textAR{نَشَآءُ} \textAR{لَقُلْنَا} \textAR{مِثْلَ} \textAR{هَـذَآ} ("We have heard (the Qur'an); if we wish we can say the like of this.") They boasted with their words, but not with their actions. They were challenged several times to bring even one chapter like the Qur'an, and they had no way to meet this challenge. They only boasted in order to deceive themselves and those who followed their falsehood. It was said that An-Nadr bin Al-Harith, may Allah curse him, was the one who said this, according to Sa`id bin Jubayr, As-Suddi, Ibn Jurayj and others. An-Nadr visited Persia and learned the stories of some Persian kings, such as Rustum and Isphandiyar. When he went back to Makkah, He found that the Prophet was sent from Allah and reciting the Qur'an to the people. Whenever the Prophet would leave an audience in which An-Nadr was sitting, An-Nadr began narrating to them the stories that he learned in Persia, proclaiming afterwards, "Who, by Allah, has better tales to narrate, I or Muhammad" When Allah allowed the Muslims to capture An-Nadr in Badr, the Messenger of Allah (saw) commanded that his head be cut off before him, and that was done, all thanks are due to Allah. The meaning of, \textAR{أَسَـطِيرُ} \textAR{الاٌّوَّلِينَ} (. ..tales of the ancients) meaning that the Prophet has plagiarized and learned books of ancient people, and this is what he narrated to people, as they claimed. This is the pure falsehood that Allah mentioned in another Ayah, \textAR{وَقَالُواْ} \textAR{أَسَـطِيرُ} \textAR{الاٌّوَّلِينَ} \textAR{اكْتَتَبَهَا} \textAR{فَهِىَ} \textAR{تُمْلَى} \textAR{عَلَيْهِ} \textAR{بُكْرَةً} \textAR{وَأَصِيلاً} - \textAR{قُلْ} \textAR{أَنزَلَهُ} \textAR{الَّذِى} \textAR{يَعْلَمُ} \textAR{السِّرَّ} \textAR{فِى} \textAR{السَّمَـوَتِ} \textAR{وَالاٌّرْضِ} \textAR{إِنَّهُ} \textAR{كَانَ} \textAR{غَفُوراً} \textAR{رَّحِيماً} (And they say: "Tales of the ancients, which he has written down:, and they are dictated to him morning and afternoon." Say: "It (this Qur'an) has been sent down by Him (Allah) Who knows the secret of the heavens and the earth. Truly, He is Oft-Forgiving, Most Merciful.") 25:5-6 for those who repent and return to Him, He accepts repentance from them and forgives them. The Idolators ask for Allah's Judgment and Torment! Allah said, \textAR{وَإِذْ} \textAR{قَالُواْ} \textAR{اللَّهُمَّ} \textAR{إِن} \textAR{كَانَ} \textAR{هَـذَا} \textAR{هُوَ} \textAR{الْحَقَّ} \textAR{مِنْ} \textAR{عِندِكَ} \textAR{فَأَمْطِرْ} \textAR{عَلَيْنَا} \textAR{حِجَارَةً} \textAR{مِّنَ} \textAR{السَّمَآءِ} \textAR{أَوِ} \textAR{ائْتِنَا} \textAR{بِعَذَابٍ} \textAR{أَلِيمٍ} (And (remember) when they said: "O Allah! If this (the Qur'an) is indeed the truth (revealed) from You, then rain down stones on us from the sky or bring on us a painful torment.") This is indicative of the pagans' enormous ignorance, denial, stubbornness and transgression. They should have said, "O Allah! If this is the truth from You, then guide us to it and help us follow it." However, they brought Allah's judgment on themselves and asked for His punishment. Allah said in other Ayat, \textAR{وَيَسْتَعْجِلُونَكَ} \textAR{بِالْعَذَابِ} \textAR{وَلَوْلاَ} \textAR{أَجَلٌ} \textAR{مُّسَمًّى} \textAR{لَّجَآءَهُمُ} \textAR{الْعَذَابُ} \textAR{وَلَيَأْتِيَنَّهُمْ} \textAR{بَغْتَةً} \textAR{وَهُمْ} \textAR{لاَ} \textAR{يَشْعُرُونَ} (And they ask you to hasten on the torment (for them), and had it not been for a term appointed, the torment would certainly have come to them. And surely, it will come upon them suddenly while they perceive not!) 29:53, \textAR{وَقَالُواْ} \textAR{رَبَّنَا} \textAR{عَجِّل} \textAR{لَّنَا} \textAR{قِطَّنَا} \textAR{قَبْلَ} \textAR{يَوْمِ} \textAR{الْحِسَابِ} (They say: "Our Lord! Hasten to us Qittana (our record of good and bad deeds so that we may see it) before the Day of Reckoning!") 38:16, and, \textAR{سَأَلَ} \textAR{سَآئِلٌ} \textAR{بِعَذَابٍ} \textAR{وَاقِعٍ} - \textAR{لِّلْكَـفِرِينَ} \textAR{لَيْسَ} \textAR{لَهُ} \textAR{دَافِعٌ} - \textAR{مِّنَ} \textAR{اللَّهِ} \textAR{ذِي} \textAR{الْمَعَارِجِ} (A questioner asked concerning a torment about to befall. Upon the disbelievers, which none can avert. From Allah, the Lord of the ways of ascent.) 70:1-3 The ignorant ones in ancient times said similar things. The people of Shu`ayb said to him, \textAR{فَأَسْقِطْ}  [...] [Tafsir continues in full corpus]
\end{tafsirbox}


\begin{muyassarbox}[Tafsir Muyassar]
\textAR{وإذا تتلى على هؤلاء الذين كفروا بالله آيات القرآن العزيز قالوا جهلا منهم وعنادًا للحق: قد سمعنا هذا من قبل، لو نشاء لقلنا مثل هذا القرآن، ما هذا القرآن الذي تتلوه علينا -أيها الرسول- إلا أكاذيب الأولين.}
\end{muyassarbox}


\newpage

\section{Verse 8:32}
{\small\color{accent2}\itshape Theme 4: Quraysh's Plots and Allah's Protection (8:30--33)}


\begin{quranbox}[Qur'an 8:32]
\quranverse{وَإِذْ قَالُوا۟ ٱللَّهُمَّ إِن كَانَ هَـٰذَا هُوَ ٱلْحَقَّ مِنْ عِندِكَ فَأَمْطِرْ عَلَيْنَا حِجَارَةً مِّنَ ٱلسَّمَآءِ أَوِ ٱئْتِنَا بِعَذَابٍ أَلِيمٍ}

\vspace{6pt}
\noindent\textbf{Sahih International:} And [remember] when they said, "O Allah, if this should be the truth from You, then rain down upon us stones from the sky or bring us a painful punishment."

\vspace{4pt}
\noindent\textbf{Pickthall:} And when they said: O Allah! If this be indeed the truth from Thee, then rain down stones on us or bring on us some painful doom!

\vspace{4pt}
\noindent\textbf{Yusuf Ali:} Remember how they said: "O Allah if this is indeed the Truth from Thee, rain down on us a shower of stones form the sky, or send us a grievous penalty."

\vspace{4pt}
\noindent\textbf{Asad:} And, lo, they would say, "O God! If this. be indeed the truth from Thee, then rain down upon us stones from the skies, or inflict [some other] grievous suffering on us!"

\vspace{4pt}
\noindent\textbf{Basmeih (Malay):} Dan (ingatlah) ketika mereka (kaum musyrik Makkah) berkata: "Wahai tuhan kami! Jika betul (Al-Quran) itu ialah yang benar dari sisimu, maka hujanilah kami dengan batu dari langit, atau datangkanlah kepada kami azab seksa yang tidak terperi sakitnya".
\end{quranbox}


\begin{tafsirbox}[Tafsir Ibn Kathir (Abridged)]
The Quraysh claimed They can produce Something similar to the Qur'an Allah describes the disbelief, transgression, rebellion, as well as misguided statements that the pagans of Quraysh used to utter when they heard Allah's Ayat being recited to them, \textAR{قَدْ} \textAR{سَمِعْنَا} \textAR{لَوْ} \textAR{نَشَآءُ} \textAR{لَقُلْنَا} \textAR{مِثْلَ} \textAR{هَـذَآ} ("We have heard (the Qur'an); if we wish we can say the like of this.") They boasted with their words, but not with their actions. They were challenged several times to bring even one chapter like the Qur'an, and they had no way to meet this challenge. They only boasted in order to deceive themselves and those who followed their falsehood. It was said that An-Nadr bin Al-Harith, may Allah curse him, was the one who said this, according to Sa`id bin Jubayr, As-Suddi, Ibn Jurayj and others. An-Nadr visited Persia and learned the stories of some Persian kings, such as Rustum and Isphandiyar. When he went back to Makkah, He found that the Prophet was sent from Allah and reciting the Qur'an to the people. Whenever the Prophet would leave an audience in which An-Nadr was sitting, An-Nadr began narrating to them the stories that he learned in Persia, proclaiming afterwards, "Who, by Allah, has better tales to narrate, I or Muhammad" When Allah allowed the Muslims to capture An-Nadr in Badr, the Messenger of Allah (saw) commanded that his head be cut off before him, and that was done, all thanks are due to Allah. The meaning of, \textAR{أَسَـطِيرُ} \textAR{الاٌّوَّلِينَ} (. ..tales of the ancients) meaning that the Prophet has plagiarized and learned books of ancient people, and this is what he narrated to people, as they claimed. This is the pure falsehood that Allah mentioned in another Ayah, \textAR{وَقَالُواْ} \textAR{أَسَـطِيرُ} \textAR{الاٌّوَّلِينَ} \textAR{اكْتَتَبَهَا} \textAR{فَهِىَ} \textAR{تُمْلَى} \textAR{عَلَيْهِ} \textAR{بُكْرَةً} \textAR{وَأَصِيلاً} - \textAR{قُلْ} \textAR{أَنزَلَهُ} \textAR{الَّذِى} \textAR{يَعْلَمُ} \textAR{السِّرَّ} \textAR{فِى} \textAR{السَّمَـوَتِ} \textAR{وَالاٌّرْضِ} \textAR{إِنَّهُ} \textAR{كَانَ} \textAR{غَفُوراً} \textAR{رَّحِيماً} (And they say: "Tales of the ancients, which he has written down:, and they are dictated to him morning and afternoon." Say: "It (this Qur'an) has been sent down by Him (Allah) Who knows the secret of the heavens and the earth. Truly, He is Oft-Forgiving, Most Merciful.") 25:5-6 for those who repent and return to Him, He accepts repentance from them and forgives them. The Idolators ask for Allah's Judgment and Torment! Allah said, \textAR{وَإِذْ} \textAR{قَالُواْ} \textAR{اللَّهُمَّ} \textAR{إِن} \textAR{كَانَ} \textAR{هَـذَا} \textAR{هُوَ} \textAR{الْحَقَّ} \textAR{مِنْ} \textAR{عِندِكَ} \textAR{فَأَمْطِرْ} \textAR{عَلَيْنَا} \textAR{حِجَارَةً} \textAR{مِّنَ} \textAR{السَّمَآءِ} \textAR{أَوِ} \textAR{ائْتِنَا} \textAR{بِعَذَابٍ} \textAR{أَلِيمٍ} (And (remember) when they said: "O Allah! If this (the Qur'an) is indeed the truth (revealed) from You, then rain down stones on us from the sky or bring on us a painful torment.") This is indicative of the pagans' enormous ignorance, denial, stubbornness and transgression. They should have said, "O Allah! If this is the truth from You, then guide us to it and help us follow it." However, they brought Allah's judgment on themselves and asked for His punishment. Allah said in other Ayat, \textAR{وَيَسْتَعْجِلُونَكَ} \textAR{بِالْعَذَابِ} \textAR{وَلَوْلاَ} \textAR{أَجَلٌ} \textAR{مُّسَمًّى} \textAR{لَّجَآءَهُمُ} \textAR{الْعَذَابُ} \textAR{وَلَيَأْتِيَنَّهُمْ} \textAR{بَغْتَةً} \textAR{وَهُمْ} \textAR{لاَ} \textAR{يَشْعُرُونَ} (And they ask you to hasten on the torment (for them), and had it not been for a term appointed, the torment would certainly have come to them. And surely, it will come upon them suddenly while they perceive not!) 29:53, \textAR{وَقَالُواْ} \textAR{رَبَّنَا} \textAR{عَجِّل} \textAR{لَّنَا} \textAR{قِطَّنَا} \textAR{قَبْلَ} \textAR{يَوْمِ} \textAR{الْحِسَابِ} (They say: "Our Lord! Hasten to us Qittana (our record of good and bad deeds so that we may see it) before the Day of Reckoning!") 38:16, and, \textAR{سَأَلَ} \textAR{سَآئِلٌ} \textAR{بِعَذَابٍ} \textAR{وَاقِعٍ} - \textAR{لِّلْكَـفِرِينَ} \textAR{لَيْسَ} \textAR{لَهُ} \textAR{دَافِعٌ} - \textAR{مِّنَ} \textAR{اللَّهِ} \textAR{ذِي} \textAR{الْمَعَارِجِ} (A questioner asked concerning a torment about to befall. Upon the disbelievers, which none can avert. From Allah, the Lord of the ways of ascent.) 70:1-3 The ignorant ones in ancient times said similar things. The people of Shu`ayb said to him, \textAR{فَأَسْقِطْ}  [...] [Tafsir continues in full corpus]
\end{tafsirbox}


\begin{muyassarbox}[Tafsir Muyassar]
\textAR{واذكر -أيها الرسول- قول المشركين من قومك داعين الله: إن كان ما جاء به محمد هو الحق مِن عندك فأمطر علينا حجارة من السماء، أو ائتنا بعذاب شديد موجع.}
\end{muyassarbox}


\newpage

\section{Verse 8:33}
{\small\color{accent2}\itshape Theme 4: Quraysh's Plots and Allah's Protection (8:30--33)}


\begin{quranbox}[Qur'an 8:33]
\quranverse{وَمَا كَانَ ٱللَّهُ لِيُعَذِّبَهُمْ وَأَنتَ فِيهِمْ ۚ وَمَا كَانَ ٱللَّهُ مُعَذِّبَهُمْ وَهُمْ يَسْتَغْفِرُونَ}

\vspace{6pt}
\noindent\textbf{Sahih International:} But Allah would not punish them while you, [O Muhammad], are among them, and Allah would not punish them while they seek forgiveness.

\vspace{4pt}
\noindent\textbf{Pickthall:} But Allah would not punish them while thou wast with them, nor will He punish them while they seek forgiveness.

\vspace{4pt}
\noindent\textbf{Yusuf Ali:} But Allah was not going to send them a penalty whilst thou wast amongst them; nor was He going to send it whilst they could ask for pardon.

\vspace{4pt}
\noindent\textbf{Asad:} But God did not choose thus to chastise them whip thou [O Prophet] wert still among them, nor would God chastise them when they [might yet] ask for forgiveness.

\vspace{4pt}
\noindent\textbf{Basmeih (Malay):} Dan Allah tidak sekali-kali akan menyeksa mereka, sedang engkau (wahai Muhammad) ada di antara mereka; dan Allah tidak akan menyeksa mereka sedang mereka beristighfar (meminta ampun).
\end{quranbox}


\begin{tafsirbox}[Tafsir Ibn Kathir (Abridged)]
The Quraysh claimed They can produce Something similar to the Qur'an Allah describes the disbelief, transgression, rebellion, as well as misguided statements that the pagans of Quraysh used to utter when they heard Allah's Ayat being recited to them, \textAR{قَدْ} \textAR{سَمِعْنَا} \textAR{لَوْ} \textAR{نَشَآءُ} \textAR{لَقُلْنَا} \textAR{مِثْلَ} \textAR{هَـذَآ} ("We have heard (the Qur'an); if we wish we can say the like of this.") They boasted with their words, but not with their actions. They were challenged several times to bring even one chapter like the Qur'an, and they had no way to meet this challenge. They only boasted in order to deceive themselves and those who followed their falsehood. It was said that An-Nadr bin Al-Harith, may Allah curse him, was the one who said this, according to Sa`id bin Jubayr, As-Suddi, Ibn Jurayj and others. An-Nadr visited Persia and learned the stories of some Persian kings, such as Rustum and Isphandiyar. When he went back to Makkah, He found that the Prophet was sent from Allah and reciting the Qur'an to the people. Whenever the Prophet would leave an audience in which An-Nadr was sitting, An-Nadr began narrating to them the stories that he learned in Persia, proclaiming afterwards, "Who, by Allah, has better tales to narrate, I or Muhammad" When Allah allowed the Muslims to capture An-Nadr in Badr, the Messenger of Allah (saw) commanded that his head be cut off before him, and that was done, all thanks are due to Allah. The meaning of, \textAR{أَسَـطِيرُ} \textAR{الاٌّوَّلِينَ} (. ..tales of the ancients) meaning that the Prophet has plagiarized and learned books of ancient people, and this is what he narrated to people, as they claimed. This is the pure falsehood that Allah mentioned in another Ayah, \textAR{وَقَالُواْ} \textAR{أَسَـطِيرُ} \textAR{الاٌّوَّلِينَ} \textAR{اكْتَتَبَهَا} \textAR{فَهِىَ} \textAR{تُمْلَى} \textAR{عَلَيْهِ} \textAR{بُكْرَةً} \textAR{وَأَصِيلاً} - \textAR{قُلْ} \textAR{أَنزَلَهُ} \textAR{الَّذِى} \textAR{يَعْلَمُ} \textAR{السِّرَّ} \textAR{فِى} \textAR{السَّمَـوَتِ} \textAR{وَالاٌّرْضِ} \textAR{إِنَّهُ} \textAR{كَانَ} \textAR{غَفُوراً} \textAR{رَّحِيماً} (And they say: "Tales of the ancients, which he has written down:, and they are dictated to him morning and afternoon." Say: "It (this Qur'an) has been sent down by Him (Allah) Who knows the secret of the heavens and the earth. Truly, He is Oft-Forgiving, Most Merciful.") 25:5-6 for those who repent and return to Him, He accepts repentance from them and forgives them. The Idolators ask for Allah's Judgment and Torment! Allah said, \textAR{وَإِذْ} \textAR{قَالُواْ} \textAR{اللَّهُمَّ} \textAR{إِن} \textAR{كَانَ} \textAR{هَـذَا} \textAR{هُوَ} \textAR{الْحَقَّ} \textAR{مِنْ} \textAR{عِندِكَ} \textAR{فَأَمْطِرْ} \textAR{عَلَيْنَا} \textAR{حِجَارَةً} \textAR{مِّنَ} \textAR{السَّمَآءِ} \textAR{أَوِ} \textAR{ائْتِنَا} \textAR{بِعَذَابٍ} \textAR{أَلِيمٍ} (And (remember) when they said: "O Allah! If this (the Qur'an) is indeed the truth (revealed) from You, then rain down stones on us from the sky or bring on us a painful torment.") This is indicative of the pagans' enormous ignorance, denial, stubbornness and transgression. They should have said, "O Allah! If this is the truth from You, then guide us to it and help us follow it." However, they brought Allah's judgment on themselves and asked for His punishment. Allah said in other Ayat, \textAR{وَيَسْتَعْجِلُونَكَ} \textAR{بِالْعَذَابِ} \textAR{وَلَوْلاَ} \textAR{أَجَلٌ} \textAR{مُّسَمًّى} \textAR{لَّجَآءَهُمُ} \textAR{الْعَذَابُ} \textAR{وَلَيَأْتِيَنَّهُمْ} \textAR{بَغْتَةً} \textAR{وَهُمْ} \textAR{لاَ} \textAR{يَشْعُرُونَ} (And they ask you to hasten on the torment (for them), and had it not been for a term appointed, the torment would certainly have come to them. And surely, it will come upon them suddenly while they perceive not!) 29:53, \textAR{وَقَالُواْ} \textAR{رَبَّنَا} \textAR{عَجِّل} \textAR{لَّنَا} \textAR{قِطَّنَا} \textAR{قَبْلَ} \textAR{يَوْمِ} \textAR{الْحِسَابِ} (They say: "Our Lord! Hasten to us Qittana (our record of good and bad deeds so that we may see it) before the Day of Reckoning!") 38:16, and, \textAR{سَأَلَ} \textAR{سَآئِلٌ} \textAR{بِعَذَابٍ} \textAR{وَاقِعٍ} - \textAR{لِّلْكَـفِرِينَ} \textAR{لَيْسَ} \textAR{لَهُ} \textAR{دَافِعٌ} - \textAR{مِّنَ} \textAR{اللَّهِ} \textAR{ذِي} \textAR{الْمَعَارِجِ} (A questioner asked concerning a torment about to befall. Upon the disbelievers, which none can avert. From Allah, the Lord of the ways of ascent.) 70:1-3 The ignorant ones in ancient times said similar things. The people of Shu`ayb said to him, \textAR{فَأَسْقِطْ}  [...] [Tafsir continues in full corpus]
\end{tafsirbox}


\begin{muyassarbox}[Tafsir Muyassar]
\textAR{وما كان الله سبحانه وتعالى ليعذِّب هؤلاء المشركين، وأنت -أيها الرسول- بين ظهرانَيْهم، وما كان الله معذِّبهم، وهم يستغفرون من ذنوبهم.}
\end{muyassarbox}


\newpage

\section{Verse 8:34}
{\small\color{accent2}\itshape Theme 5: Deserved Punishment and Futile Spending (8:34--37)}


\begin{quranbox}[Qur'an 8:34]
\quranverse{وَمَا لَهُمْ أَلَّا يُعَذِّبَهُمُ ٱللَّهُ وَهُمْ يَصُدُّونَ عَنِ ٱلْمَسْجِدِ ٱلْحَرَامِ وَمَا كَانُوٓا۟ أَوْلِيَآءَهُۥٓ ۚ إِنْ أَوْلِيَآؤُهُۥٓ إِلَّا ٱلْمُتَّقُونَ وَلَـٰكِنَّ أَكْثَرَهُمْ لَا يَعْلَمُونَ}

\vspace{6pt}
\noindent\textbf{Sahih International:} But why should Allah not punish them while they obstruct [people] from al-Masjid al- Haram and they were not [fit to be] its guardians? Its [true] guardians are not but the righteous, but most of them do not know.

\vspace{4pt}
\noindent\textbf{Pickthall:} What (plea) have they that Allah should not punish them, when they debar (His servants) from the Inviolable Place of Worship, though they are not its fitting guardians. Its fitting guardians are those only who keep their duty to Allah. But most of them know not.

\vspace{4pt}
\noindent\textbf{Yusuf Ali:} But what plea have they that Allah should not punish them, when they keep out (men) from the sacred Mosque - and they are not its guardians? No men can be its guardians except the righteous; but most of them do not understand.

\vspace{4pt}
\noindent\textbf{Asad:} But what have they [now] in their favour that God should not chastise them-seeing that they bar [the believers] from the Inviolable House of Worship, although they are not its [rightful] guardians? None but the God-conscious can be its guardians: but of this most of these [evildoers] are unaware;

\vspace{4pt}
\noindent\textbf{Basmeih (Malay):} Dan mengapa mereka tidak patut diseksa oleh Allah, sedang mereka menyekat (orang-orang Islam) dari masjid Al-Haraam, padahal mereka bukanlah orang-orang yang berhak menguasainya (kerana mereka kafir musyrik)? Sebenarnya orang-orang yang berhak menguasainya hanyalah orang-orang yang bertaqwa, tetapi kebanyakan mereka tidak mengetahui.
\end{quranbox}


\begin{tafsirbox}[Tafsir Ibn Kathir (Abridged)]
The Idolators deserved Allah's Torment after Their Atrocities Allah states that the idolators deserved the torment, but He did not torment them in honor of the Prophet residing among them. After Allah allowed the Prophet to migrate away from them, He sent His torment upon them on the day of Badr. During that battle, the chief pagans were killed, or captured. Allah also directed them to seek forgiveness for the sins, Shirk and wickedness they indulged in. If it was not for the fact that there were some weak Muslims living among the Makkan pagans, those Muslims who invoked Allah for His forgiveness, Allah would have sent down to them the torment that could never be averted. Allah did not do that on account of the weak, ill-treated, and oppressed believers living among them, as He reiterated about the day at Al-Hudaybiyyah, \textAR{هُمُ} \textAR{الَّذِينَ} \textAR{كَفَرُواْ} \textAR{وَصَدُّوكُمْ} \textAR{عَنِ} \textAR{الْمَسْجِدِ} \textAR{الْحَرَامِ} \textAR{وَالْهَدْىَ} \textAR{مَعْكُوفاً} \textAR{أَن} \textAR{يَبْلُغَ} \textAR{مَحِلَّهُ} \textAR{وَلَوْلاَ} \textAR{رِجَالٌ} \textAR{مُّؤْمِنُونَ} \textAR{وَنِسَآءٌ} \textAR{مُّؤْمِنَـتٌ} \textAR{لَّمْ} \textAR{تَعْلَمُوهُمْ} \textAR{أَن} \textAR{تَطَئُوهُمْ} \textAR{فَتُصِيبَكمْ} \textAR{مِّنْهُمْ} \textAR{مَّعَرَّةٌ} \textAR{بِغَيْرِ} \textAR{عِلْمٍ} \textAR{لِّيُدْخِلَ} \textAR{اللَّهُ} \textAR{فِى} \textAR{رَحْمَتِهِ} \textAR{مَن} \textAR{يَشَآءُ} \textAR{لَوْ} \textAR{تَزَيَّلُواْ} \textAR{لَعَذَّبْنَا} \textAR{الَّذِينَ} \textAR{كَفَرُواْ} \textAR{مِنْهُمْ} \textAR{عَذَاباً} \textAR{أَلِيماً} (They are the ones who disbelieved and hindered you from Al-Masjid Al-Haram (at Makkah) and detained the sacrificial animals from reaching their place of sacrifice. Had there not been believing men and believing women whom you did not know, that you may kill them and on whose account a sin would have been committed by you without (your) knowledge, that Allah might bring into His mercy whom He wills if they (the believers and the disbelievers) had been apart, We verily, would have punished those of them who disbelieved with painful torment.) 48:25 Allah said here, \textAR{وَمَا} \textAR{لَهُمْ} \textAR{أَلاَّ} \textAR{يُعَذِّبَهُمُ} \textAR{اللَّهُ} \textAR{وَهُمْ} \textAR{يَصُدُّونَ} \textAR{عَنِ} \textAR{الْمَسْجِدِ} \textAR{الْحَرَامِ} \textAR{وَمَا} \textAR{كَانُواْ} \textAR{أَوْلِيَآءَهُ} \textAR{إِنْ} \textAR{أَوْلِيَآؤُهُ} \textAR{إِلاَّ} \textAR{الْمُتَّقُونَ} \textAR{وَلَـكِنَّ} \textAR{أَكْثَرَهُمْ} \textAR{لاَ} \textAR{يَعْلَمُونَ} (And why should not Allah punish them while they hinder (men) from Al-Masjid Al-Haram, and they are not its guardians None can be its guardians except those who have Taqwa, but most of them know not.) Allah asks, `why would not He torment them while they are stopping Muslims from going to Al-Masjid Al-Haram, thus hindering the believers, its own people, from praying and performing Tawaf in it' Allah said, \textAR{وَمَا} \textAR{كَانُواْ} \textAR{أَوْلِيَآءَهُ} \textAR{إِنْ} \textAR{أَوْلِيَآؤُهُ} \textAR{إِلاَّ} \textAR{الْمُتَّقُونَ} T(And they are not its guardians None can be its guardians except those who have Taqwa,) meaning, the Prophet and his Companions are the true dwellers (or worthy maintainers) of Al-Masjid Al-Haram, not the pagans. Allah said in other Ayah, \textAR{مَا} \textAR{كَانَ} \textAR{لِلْمُشْرِكِينَ} \textAR{أَن} \textAR{يَعْمُرُواْ} \textAR{مَسَاجِدَ} \textAR{الله} \textAR{شَـهِدِينَ} \textAR{عَلَى} \textAR{أَنفُسِهِم} \textAR{بِالْكُفْرِ} \textAR{أُوْلَـئِكَ} \textAR{حَبِطَتْ} \textAR{أَعْمَـلُهُمْ} \textAR{وَفِى} \textAR{النَّارِ} \textAR{هُمْ} \textAR{خَـلِدُونَ} - \textAR{إِنَّمَا} \textAR{يَعْمُرُ} \textAR{مَسَـجِدَ} \textAR{اللَّهِ} \textAR{مَنْ} \textAR{ءَامَنَ} \textAR{بِاللَّهِ} \textAR{وَالْيَوْمِ} \textAR{الاٌّخِرِ} \textAR{وَأَقَامَ} \textAR{الصَّلَوةَ} \textAR{وَءاتَى} \textAR{الزَّكَوةَ} \textAR{وَلَمْ} \textAR{يَخْشَ} \textAR{إِلاَّ} \textAR{اللَّهَ} \textAR{فَعَسَى} \textAR{أُوْلَـئِكَ} \textAR{أَن} \textAR{يَكُونُواْ} \textAR{مِنَ} \textAR{الْمُهْتَدِينَ} (It is not for the polytheists, to maintain the Masjids of Allah, while they witness disbelief against themselves. The works of such are in vain and in the Fire shall they abide. The Masjids of Allah shall be maintained only by those who believe in Allah and the Last Day; perform the Salah, and give the Zakah and fear none but Allah. It is they who are on true guidance.) 9:17-18, and, \textAR{وَصَدٌّ} \textAR{عَن} \textAR{سَبِيلِ} \textAR{اللَّهِ} \textAR{وَكُفْرٌ} \textAR{بِهِ} \textAR{وَالْمَسْجِدِ} \textAR{الْحَرَامِ} \textAR{وَإِخْرَاجُ} \textAR{أَهْلِهِ} \textAR{مِنْهُ} \textAR{أَكْبَرُ} \textAR{عِندَ} \textAR{اللَّهِ} (But a greater (transgression) with Allah is to prevent mankind from following the way of Allah, to disbelieve in Him, to prevent access to Al-Masjid Al-Haram (at Makkah), and to drive out its inhabitants,) 2:217. `Urwah, As-Suddi and Muhammad bin Ishaq said that Allah's statement, \textAR{إِنْ} \textAR{أَوْلِيَآؤُهُ} \textAR{إِلاَّ} \textAR{الْمُتَّقُونَ} (None  [...] [Tafsir continues in full corpus]
\end{tafsirbox}


\begin{muyassarbox}[Tafsir Muyassar]
\textAR{وكيف لا يستحقُّون عذاب الله، وهم يصدون أولياءه المؤمنين عن الطواف بالكعبة والصلاة في المسجد الحرام؟ وما كانوا أولياء الله، إنْ أولياء الله إلا الذين يتقونه بأداء فرائضه واجتناب معاصيه، ولكن أكثر الكفار لا يعلمون; فلذلك ادَّعوا لأنفسهم أمرًا، غيرهم أولى به.}
\end{muyassarbox}


\newpage

\section{Verse 8:35}
{\small\color{accent2}\itshape Theme 5: Deserved Punishment and Futile Spending (8:34--37)}


\begin{quranbox}[Qur'an 8:35]
\quranverse{وَمَا كَانَ صَلَاتُهُمْ عِندَ ٱلْبَيْتِ إِلَّا مُكَآءً وَتَصْدِيَةً ۚ فَذُوقُوا۟ ٱلْعَذَابَ بِمَا كُنتُمْ تَكْفُرُونَ}

\vspace{6pt}
\noindent\textbf{Sahih International:} And their prayer at the House was not except whistling and handclapping. So taste the punishment for what you disbelieved.

\vspace{4pt}
\noindent\textbf{Pickthall:} And their worship at the (holy) House is naught but whistling and hand-clapping. Therefore (it is said unto them): Taste of the doom because ye disbelieve.

\vspace{4pt}
\noindent\textbf{Yusuf Ali:} Their prayer at the House (of Allah) is nothing but whistling and clapping of hands: (Its only answer can be), "Taste ye the penalty because ye blasphemed."

\vspace{4pt}
\noindent\textbf{Asad:} and their prayers before the Temple are nothing but whistling and clapping of hands. Taste then, [O unbelievers,] this chastisement as an outcome of your persistent denial of the truth!

\vspace{4pt}
\noindent\textbf{Basmeih (Malay):} Dan tiadalah sembahyang mereka di sisi Baitullah itu melainkan bersiul-siul dan bertepuk tangan. Oleh itu rasalah kamu (wahai orang kafir) akan azab seksa dengan sebab kekufuran kamu.
\end{quranbox}


\begin{tafsirbox}[Tafsir Ibn Kathir (Abridged)]
The Idolators deserved Allah's Torment after Their Atrocities Allah states that the idolators deserved the torment, but He did not torment them in honor of the Prophet residing among them. After Allah allowed the Prophet to migrate away from them, He sent His torment upon them on the day of Badr. During that battle, the chief pagans were killed, or captured. Allah also directed them to seek forgiveness for the sins, Shirk and wickedness they indulged in. If it was not for the fact that there were some weak Muslims living among the Makkan pagans, those Muslims who invoked Allah for His forgiveness, Allah would have sent down to them the torment that could never be averted. Allah did not do that on account of the weak, ill-treated, and oppressed believers living among them, as He reiterated about the day at Al-Hudaybiyyah, \textAR{هُمُ} \textAR{الَّذِينَ} \textAR{كَفَرُواْ} \textAR{وَصَدُّوكُمْ} \textAR{عَنِ} \textAR{الْمَسْجِدِ} \textAR{الْحَرَامِ} \textAR{وَالْهَدْىَ} \textAR{مَعْكُوفاً} \textAR{أَن} \textAR{يَبْلُغَ} \textAR{مَحِلَّهُ} \textAR{وَلَوْلاَ} \textAR{رِجَالٌ} \textAR{مُّؤْمِنُونَ} \textAR{وَنِسَآءٌ} \textAR{مُّؤْمِنَـتٌ} \textAR{لَّمْ} \textAR{تَعْلَمُوهُمْ} \textAR{أَن} \textAR{تَطَئُوهُمْ} \textAR{فَتُصِيبَكمْ} \textAR{مِّنْهُمْ} \textAR{مَّعَرَّةٌ} \textAR{بِغَيْرِ} \textAR{عِلْمٍ} \textAR{لِّيُدْخِلَ} \textAR{اللَّهُ} \textAR{فِى} \textAR{رَحْمَتِهِ} \textAR{مَن} \textAR{يَشَآءُ} \textAR{لَوْ} \textAR{تَزَيَّلُواْ} \textAR{لَعَذَّبْنَا} \textAR{الَّذِينَ} \textAR{كَفَرُواْ} \textAR{مِنْهُمْ} \textAR{عَذَاباً} \textAR{أَلِيماً} (They are the ones who disbelieved and hindered you from Al-Masjid Al-Haram (at Makkah) and detained the sacrificial animals from reaching their place of sacrifice. Had there not been believing men and believing women whom you did not know, that you may kill them and on whose account a sin would have been committed by you without (your) knowledge, that Allah might bring into His mercy whom He wills if they (the believers and the disbelievers) had been apart, We verily, would have punished those of them who disbelieved with painful torment.) 48:25 Allah said here, \textAR{وَمَا} \textAR{لَهُمْ} \textAR{أَلاَّ} \textAR{يُعَذِّبَهُمُ} \textAR{اللَّهُ} \textAR{وَهُمْ} \textAR{يَصُدُّونَ} \textAR{عَنِ} \textAR{الْمَسْجِدِ} \textAR{الْحَرَامِ} \textAR{وَمَا} \textAR{كَانُواْ} \textAR{أَوْلِيَآءَهُ} \textAR{إِنْ} \textAR{أَوْلِيَآؤُهُ} \textAR{إِلاَّ} \textAR{الْمُتَّقُونَ} \textAR{وَلَـكِنَّ} \textAR{أَكْثَرَهُمْ} \textAR{لاَ} \textAR{يَعْلَمُونَ} (And why should not Allah punish them while they hinder (men) from Al-Masjid Al-Haram, and they are not its guardians None can be its guardians except those who have Taqwa, but most of them know not.) Allah asks, `why would not He torment them while they are stopping Muslims from going to Al-Masjid Al-Haram, thus hindering the believers, its own people, from praying and performing Tawaf in it' Allah said, \textAR{وَمَا} \textAR{كَانُواْ} \textAR{أَوْلِيَآءَهُ} \textAR{إِنْ} \textAR{أَوْلِيَآؤُهُ} \textAR{إِلاَّ} \textAR{الْمُتَّقُونَ} T(And they are not its guardians None can be its guardians except those who have Taqwa,) meaning, the Prophet and his Companions are the true dwellers (or worthy maintainers) of Al-Masjid Al-Haram, not the pagans. Allah said in other Ayah, \textAR{مَا} \textAR{كَانَ} \textAR{لِلْمُشْرِكِينَ} \textAR{أَن} \textAR{يَعْمُرُواْ} \textAR{مَسَاجِدَ} \textAR{الله} \textAR{شَـهِدِينَ} \textAR{عَلَى} \textAR{أَنفُسِهِم} \textAR{بِالْكُفْرِ} \textAR{أُوْلَـئِكَ} \textAR{حَبِطَتْ} \textAR{أَعْمَـلُهُمْ} \textAR{وَفِى} \textAR{النَّارِ} \textAR{هُمْ} \textAR{خَـلِدُونَ} - \textAR{إِنَّمَا} \textAR{يَعْمُرُ} \textAR{مَسَـجِدَ} \textAR{اللَّهِ} \textAR{مَنْ} \textAR{ءَامَنَ} \textAR{بِاللَّهِ} \textAR{وَالْيَوْمِ} \textAR{الاٌّخِرِ} \textAR{وَأَقَامَ} \textAR{الصَّلَوةَ} \textAR{وَءاتَى} \textAR{الزَّكَوةَ} \textAR{وَلَمْ} \textAR{يَخْشَ} \textAR{إِلاَّ} \textAR{اللَّهَ} \textAR{فَعَسَى} \textAR{أُوْلَـئِكَ} \textAR{أَن} \textAR{يَكُونُواْ} \textAR{مِنَ} \textAR{الْمُهْتَدِينَ} (It is not for the polytheists, to maintain the Masjids of Allah, while they witness disbelief against themselves. The works of such are in vain and in the Fire shall they abide. The Masjids of Allah shall be maintained only by those who believe in Allah and the Last Day; perform the Salah, and give the Zakah and fear none but Allah. It is they who are on true guidance.) 9:17-18, and, \textAR{وَصَدٌّ} \textAR{عَن} \textAR{سَبِيلِ} \textAR{اللَّهِ} \textAR{وَكُفْرٌ} \textAR{بِهِ} \textAR{وَالْمَسْجِدِ} \textAR{الْحَرَامِ} \textAR{وَإِخْرَاجُ} \textAR{أَهْلِهِ} \textAR{مِنْهُ} \textAR{أَكْبَرُ} \textAR{عِندَ} \textAR{اللَّهِ} (But a greater (transgression) with Allah is to prevent mankind from following the way of Allah, to disbelieve in Him, to prevent access to Al-Masjid Al-Haram (at Makkah), and to drive out its inhabitants,) 2:217. `Urwah, As-Suddi and Muhammad bin Ishaq said that Allah's statement, \textAR{إِنْ} \textAR{أَوْلِيَآؤُهُ} \textAR{إِلاَّ} \textAR{الْمُتَّقُونَ} (None  [...] [Tafsir continues in full corpus]
\end{tafsirbox}


\begin{muyassarbox}[Tafsir Muyassar]
\textAR{وما كان صلاتهم عند المسجد الحرام إلا صفيرًا وتصفيقًا. فذوقوا عذاب القتل والأسر يوم «بدر» ; بسبب جحودكم وأفعالكم التي لا يُقْدم عليها إلا الكفرة، الجاحدون توحيد ربهم ورسالة نبيهم.}
\end{muyassarbox}


\newpage

\section{Verse 8:36}
{\small\color{accent2}\itshape Theme 5: Deserved Punishment and Futile Spending (8:34--37)}


\begin{quranbox}[Qur'an 8:36]
\quranverse{إِنَّ ٱلَّذِينَ كَفَرُوا۟ يُنفِقُونَ أَمْوَٰلَهُمْ لِيَصُدُّوا۟ عَن سَبِيلِ ٱللَّهِ ۚ فَسَيُنفِقُونَهَا ثُمَّ تَكُونُ عَلَيْهِمْ حَسْرَةً ثُمَّ يُغْلَبُونَ ۗ وَٱلَّذِينَ كَفَرُوٓا۟ إِلَىٰ جَهَنَّمَ يُحْشَرُونَ}

\vspace{6pt}
\noindent\textbf{Sahih International:} Indeed, those who disbelieve spend their wealth to avert [people] from the way of Allah. So they will spend it; then it will be for them a [source of] regret; then they will be overcome. And those who have disbelieved - unto Hell they will be gathered.

\vspace{4pt}
\noindent\textbf{Pickthall:} Lo! those who disbelieve spend their wealth in order that they may debar (men) from the way of Allah. They will spend it, then it will become an anguish for them, then they will be conquered. And those who disbelieve will be gathered unto hell,

\vspace{4pt}
\noindent\textbf{Yusuf Ali:} The Unbelievers spend their wealth to hinder (man) from the path of Allah, and so will they continue to spend; but in the end they will have (only) regrets and sighs; at length they will be overcome: and the Unbelievers will be gathered together to Hell;-

\vspace{4pt}
\noindent\textbf{Asad:} Behold, those who are bent on denying the truth are spending their riches in order to turn others away from the path of God; and they will go on spending them until they become [a source of] intense regret for them; and then they will be overcome! And those who [until their death] have denied the truth shall be gathered unto hell,

\vspace{4pt}
\noindent\textbf{Basmeih (Malay):} Sesungguhnya orang-orang kafir yang selalu membelanjakan harta mereka untuk menghalangi (manusia) dari jalan Allah, maka mereka tetap membelanjakannya kemudian (harta yang dibelanjakan) itu menyebabkan penyesalan kepada mereka, tambahan pula mereka dikalahkan. Dan (ingatlah) orang-orang kafir itu (akhirnya) dihimpunkan dalam neraka jahanam.
\end{quranbox}


\begin{tafsirbox}[Tafsir Ibn Kathir (Abridged)]
The Disbelievers spend Their Wealth to hinder Others from Allah's Path, but this will only cause Them Grief Muhammad bin Ishaq narrated that Az-Zuhri, Muhammad bin Yahya bin Hibban, `Asim bin `Umar bin Qatadah, and Al-Husayn bin `Abdur-Rahman bin `Amr bin Sa`id bin Mu`adh said, "The Quraysh suffered defeat at Badr and their forces went back to Makkah, while Abu Sufyan went back with the caravan intact. This is when `Abdullah bin Abi Rabi`ah, `Ikrimah bin Abi Jahl, Safwan bin Umayyah and other men from Quraysh who lost their fathers, sons or brothers in Badr, went to Abu Sufyan bin Harb. They said to him, and to those among the Quraysh who had wealth in that caravan, `O people of Quraysh! Muhammad has grieved you and killed the chiefs among you. Therefore, help us with this wealth so that we can fight him, it may be that we will avenge our losses.' They agreed." Muhammad bin Ishaq said, "This Ayah was revealed about them, according to Ibn `Abbas, \textAR{إِنَّ} \textAR{الَّذِينَ} \textAR{كَفَرُواْ} \textAR{يُنفِقُونَ} \textAR{أَمْوَلَهُمْ} (Verily, those who disbelieve spend their wealth...) until, \textAR{هُمُ} \textAR{الْخَـسِرُونَ} (they who are the losers. )" Mujahid, Sa`id bin Jubayr, Al-Hakam bin `Uyaynah, Qatadah, As-Suddi and Ibn Abza said that this Ayah was revealed about Abu Sufyan and his spending money in Uhud to fight the Messenger of Allah (saw). Ad-Dahhak said that this Ayah was revealed about the idolators of Badr. In any case, the Ayah is general, even though there was a specific incident that accompanied its revelation. Allah states here that the disbelievers spend their wealth to hinder from the path of truth. However, by doing that, their money will be spent and then will become a source of grief and anguish for them, availing them nothing in the least. They seek to extinguish the Light of Allah and make their word higher than the word of truth. However, Allah will complete His Light, even though the disbelievers hate it. He will give aid to His religion, make His Word dominant, and His religion will prevail above all religions. This is the disgrace that the disbelievers will taste in this life; and in the Hereafter, they will taste the torment of the Fire. Whoever among them lives long, will witness with his eyes and hear with his ears what causes grief to him. Those among them who are killed or die will be returned to eternal disgrace and the everlasting punishment. This is why Allah said, \textAR{فَسَيُنفِقُونَهَا} \textAR{ثُمَّ} \textAR{تَكُونُ} \textAR{عَلَيْهِمْ} \textAR{حَسْرَةً} \textAR{ثُمَّ} \textAR{يُغْلَبُونَ} \textAR{وَالَّذِينَ} \textAR{كَفَرُواْ} \textAR{إِلَى} \textAR{جَهَنَّمَ} \textAR{يُحْشَرُونَ} (And so will they continue to spend it; but in the end it will become an anguish for them. Then they will be overcome. And those who disbelieve will be gathered unto Hell. ) Allah said, \textAR{لِيَمِيزَ} \textAR{اللَّهُ} \textAR{الْخَبِيثَ} \textAR{مِنَ} \textAR{الطَّيِّبِ} (In order that Allah may distinguish the wicked from the good.), meaning recognize the difference between the people of happiness and the people of misery, according to Ibn `Abbas, as `Ali bin Abi Talhah reported from him. Allah distinguishes between those believers who obey Him and fight His disbelieving enemies and those who disobey Him. Allah said in another Ayah, \textAR{مَّا} \textAR{كَانَ} \textAR{اللَّهُ} \textAR{لِيَذَرَ} \textAR{الْمُؤْمِنِينَ} \textAR{عَلَى} \textAR{مَآ} \textAR{أَنتُمْ} \textAR{عَلَيْهِ} \textAR{حَتَّى} \textAR{يَمِيزَ} \textAR{الْخَبِيثَ} \textAR{مِنَ} \textAR{الطَّيِّبِ} \textAR{وَمَا} \textAR{كَانَ} \textAR{اللَّهُ} \textAR{لِيُطْلِعَكُمْ} \textAR{عَلَى} \textAR{الْغَيْبِ} (Allah will not leave the believers in the state in which you are now, until He distinguishes the wicked from the good. Nor will Allah disclose to you the secrets of the Ghayb (Unseen).) 3:179, and, \textAR{أَمْ} \textAR{حَسِبْتُمْ} \textAR{أَن} \textAR{تَدْخُلُواْ} \textAR{الْجَنَّةَ} \textAR{وَلَمَّا} \textAR{يَعْلَمِ} \textAR{اللَّهُ} \textAR{الَّذِينَ} \textAR{جَـهَدُواْ} \textAR{مِنكُمْ} \textAR{وَيَعْلَمَ} \textAR{الصَّـبِرِينَ} (Do you think that you will enter Paradise before Allah (tests) those of you who fought (in His cause) and (also) tests those who are the patient)3:142. Therefore, the Ayah (8:37) means, `We tried you with combatant disbelievers whom We made able to spend money in fighting you,' \textAR{لِيَمِيزَ} \textAR{اللَّهُ} \textAR{الْخَ} [...] [Tafsir continues in full corpus]
\end{tafsirbox}


\begin{muyassarbox}[Tafsir Muyassar]
\textAR{إن الذين جحدوا وحدانية الله وعصوا رسوله ينفقون أموالهم فيعطونها أمثالهم من المشركين وأهل الضلال، ليصدوا عن سبيل الله ويمنعوا المؤمنين عن الإيمان بالله ورسوله، فينفقون أموالهم في ذلك، ثم تكون عاقبة نفقتهم تلك ندامة وحسرة عليهم; لأن أموالهم تذهب، ولا يظفرون بما يأمُلون مِن إطفاء نور الله والصد عن سبيله، ثم يهزمهم المؤمنون آخر الأمر. والذين كفروا إلى جهنم يحشرون فيعذبون فيها.}
\end{muyassarbox}


\newpage

\section{Verse 8:37}
{\small\color{accent2}\itshape Theme 5: Deserved Punishment and Futile Spending (8:34--37)}


\begin{quranbox}[Qur'an 8:37]
\quranverse{لِيَمِيزَ ٱللَّهُ ٱلْخَبِيثَ مِنَ ٱلطَّيِّبِ وَيَجْعَلَ ٱلْخَبِيثَ بَعْضَهُۥ عَلَىٰ بَعْضٍ فَيَرْكُمَهُۥ جَمِيعًا فَيَجْعَلَهُۥ فِى جَهَنَّمَ ۚ أُو۟لَـٰٓئِكَ هُمُ ٱلْخَـٰسِرُونَ}

\vspace{6pt}
\noindent\textbf{Sahih International:} [This is] so that Allah may distinguish the wicked from the good and place the wicked some of them upon others and heap them all together and put them into Hell. It is those who are the losers.

\vspace{4pt}
\noindent\textbf{Pickthall:} That Allah may separate the wicked from the good, The wicked will He place piece upon piece, and heap them all together, and consign them unto hell. Such verily are the losers.

\vspace{4pt}
\noindent\textbf{Yusuf Ali:} In order that Allah may separate the impure from the pure, put the impure, one on another, heap them together, and cast them into Hell. They will be the ones to have lost.

\vspace{4pt}
\noindent\textbf{Asad:} so that God might separate the bad from the good, and join the bad with one another, and link them all together [within His condemnation], and then place them in hell. They, they are the lost!

\vspace{4pt}
\noindent\textbf{Basmeih (Malay):} Kerana Allah hendak membezakan yang jahat (golongan yang ingkar) dari yang baik (golongan yang beriman), dan menjadikan (golongan) yang jahat itu setengahnya bersatu dengan setengahnya yang lain, lalu ditimbunkannya kesemuanya, serta dimasukkannya ke dalam neraka Jahanam. Mereka itulah orang-orang yang rugi.
\end{quranbox}


\begin{tafsirbox}[Tafsir Ibn Kathir (Abridged)]
The Disbelievers spend Their Wealth to hinder Others from Allah's Path, but this will only cause Them Grief Muhammad bin Ishaq narrated that Az-Zuhri, Muhammad bin Yahya bin Hibban, `Asim bin `Umar bin Qatadah, and Al-Husayn bin `Abdur-Rahman bin `Amr bin Sa`id bin Mu`adh said, "The Quraysh suffered defeat at Badr and their forces went back to Makkah, while Abu Sufyan went back with the caravan intact. This is when `Abdullah bin Abi Rabi`ah, `Ikrimah bin Abi Jahl, Safwan bin Umayyah and other men from Quraysh who lost their fathers, sons or brothers in Badr, went to Abu Sufyan bin Harb. They said to him, and to those among the Quraysh who had wealth in that caravan, `O people of Quraysh! Muhammad has grieved you and killed the chiefs among you. Therefore, help us with this wealth so that we can fight him, it may be that we will avenge our losses.' They agreed." Muhammad bin Ishaq said, "This Ayah was revealed about them, according to Ibn `Abbas, \textAR{إِنَّ} \textAR{الَّذِينَ} \textAR{كَفَرُواْ} \textAR{يُنفِقُونَ} \textAR{أَمْوَلَهُمْ} (Verily, those who disbelieve spend their wealth...) until, \textAR{هُمُ} \textAR{الْخَـسِرُونَ} (they who are the losers. )" Mujahid, Sa`id bin Jubayr, Al-Hakam bin `Uyaynah, Qatadah, As-Suddi and Ibn Abza said that this Ayah was revealed about Abu Sufyan and his spending money in Uhud to fight the Messenger of Allah (saw). Ad-Dahhak said that this Ayah was revealed about the idolators of Badr. In any case, the Ayah is general, even though there was a specific incident that accompanied its revelation. Allah states here that the disbelievers spend their wealth to hinder from the path of truth. However, by doing that, their money will be spent and then will become a source of grief and anguish for them, availing them nothing in the least. They seek to extinguish the Light of Allah and make their word higher than the word of truth. However, Allah will complete His Light, even though the disbelievers hate it. He will give aid to His religion, make His Word dominant, and His religion will prevail above all religions. This is the disgrace that the disbelievers will taste in this life; and in the Hereafter, they will taste the torment of the Fire. Whoever among them lives long, will witness with his eyes and hear with his ears what causes grief to him. Those among them who are killed or die will be returned to eternal disgrace and the everlasting punishment. This is why Allah said, \textAR{فَسَيُنفِقُونَهَا} \textAR{ثُمَّ} \textAR{تَكُونُ} \textAR{عَلَيْهِمْ} \textAR{حَسْرَةً} \textAR{ثُمَّ} \textAR{يُغْلَبُونَ} \textAR{وَالَّذِينَ} \textAR{كَفَرُواْ} \textAR{إِلَى} \textAR{جَهَنَّمَ} \textAR{يُحْشَرُونَ} (And so will they continue to spend it; but in the end it will become an anguish for them. Then they will be overcome. And those who disbelieve will be gathered unto Hell. ) Allah said, \textAR{لِيَمِيزَ} \textAR{اللَّهُ} \textAR{الْخَبِيثَ} \textAR{مِنَ} \textAR{الطَّيِّبِ} (In order that Allah may distinguish the wicked from the good.), meaning recognize the difference between the people of happiness and the people of misery, according to Ibn `Abbas, as `Ali bin Abi Talhah reported from him. Allah distinguishes between those believers who obey Him and fight His disbelieving enemies and those who disobey Him. Allah said in another Ayah, \textAR{مَّا} \textAR{كَانَ} \textAR{اللَّهُ} \textAR{لِيَذَرَ} \textAR{الْمُؤْمِنِينَ} \textAR{عَلَى} \textAR{مَآ} \textAR{أَنتُمْ} \textAR{عَلَيْهِ} \textAR{حَتَّى} \textAR{يَمِيزَ} \textAR{الْخَبِيثَ} \textAR{مِنَ} \textAR{الطَّيِّبِ} \textAR{وَمَا} \textAR{كَانَ} \textAR{اللَّهُ} \textAR{لِيُطْلِعَكُمْ} \textAR{عَلَى} \textAR{الْغَيْبِ} (Allah will not leave the believers in the state in which you are now, until He distinguishes the wicked from the good. Nor will Allah disclose to you the secrets of the Ghayb (Unseen).) 3:179, and, \textAR{أَمْ} \textAR{حَسِبْتُمْ} \textAR{أَن} \textAR{تَدْخُلُواْ} \textAR{الْجَنَّةَ} \textAR{وَلَمَّا} \textAR{يَعْلَمِ} \textAR{اللَّهُ} \textAR{الَّذِينَ} \textAR{جَـهَدُواْ} \textAR{مِنكُمْ} \textAR{وَيَعْلَمَ} \textAR{الصَّـبِرِينَ} (Do you think that you will enter Paradise before Allah (tests) those of you who fought (in His cause) and (also) tests those who are the patient)3:142. Therefore, the Ayah (8:37) means, `We tried you with combatant disbelievers whom We made able to spend money in fighting you,' \textAR{لِيَمِيزَ} \textAR{اللَّهُ} \textAR{الْخَ} [...] [Tafsir continues in full corpus]
\end{tafsirbox}


\begin{muyassarbox}[Tafsir Muyassar]
\textAR{يحشر الله ويخزي هؤلاء الذين كفروا بربهم، وأنفقوا أموالهم لمنع الناس عن الإيمان بالله والصد عن سبيله; ليميز الله تعالى الخبيث من الطيب، ويجعل الله المال الحرام الذي أُنفق للصدِّ عن دين الله بعضه فوق بعض متراكمًا متراكبًا، فيجعله في نار جهنم، هؤلاء الكفار هم الخاسرون في الدنيا والآخرة.}
\end{muyassarbox}


\chapter{Thematic Synthesis and Lessons}

\section{Theme 1: The Call to Life and Warning of Fitnah (8:24--25)}
\begin{themebox}[Key Lessons]
\textbf{Respond to what gives life.} The call to respond to Allah and His Messenger is a call to spiritual vitality. Al-Bukhari interpreted ``that which gives you life'' as ``that which makes your affairs good.'' The Prophet taught Abu Sa`id that this response takes priority even over prayer.

\textbf{Allah intervenes between man and his heart.} Allah is intimately involved in the human heart. He can turn a person's heart, and the believer must recognise that the heart is not solely in human control.

\textbf{Fitnah is not selective.} A trial ``will not strike those who have wronged exclusively.'' When evil is tolerated, consequences reach the innocent. Communal responsibility: silence in the face of wrongdoing brings collective consequences.
\end{themebox}

\section{Theme 2: Divine Favour and Warning against Betrayal (8:26--28)}
\begin{themebox}[Key Lessons]
\textbf{Remember your past weakness.} The believers were few and oppressed, fearing abduction. Allah gave refuge, victory, and sustenance. Gratitude for Allah's blessings is a religious obligation.

\textbf{Do not betray your trusts.} Revealed about Hatib bin Abi Balta`ah, who warned the Quraysh of the Prophet's march. Betrayal is not merely material --- it is spiritual allegiance.

\textbf{Wealth and children are a trial.} Possessions and progeny are tests: do they draw the believer closer to Allah or become a distraction? The ``great reward'' with Allah is the real goal.
\end{themebox}

\section{Theme 3: Taqwa and the Criterion (8:29)}
\begin{themebox}[Key Lessons]
\textbf{Taqwa yields furqan.} God-consciousness produces the ability to distinguish truth from falsehood. Ibn Ishaq: furqan is the discernment from living in awareness of Allah.

\textbf{Three gifts of taqwa:} (1) Furqan --- the criterion; (2) expiation of evil deeds; (3) forgiveness. A comprehensive spiritual programme: discernment, purification, and pardon.
\end{themebox}

\section{Theme 4: The Quraysh's Plots and Allah's Protection (8:30--33)}
\begin{themebox}[Key Lessons]
\textbf{The Makkan plot backfired.} At Dar An-Nadwah, the Quraysh chose expulsion. Iblis guided this decision. The Hijra led to the founding of the Islamic state. Allah's plan turned their plot into Islam's triumph.

\textbf{The defiant prayer for punishment.} The Quraysh dared Allah to rain stones on them. An-Nadr bin Al-Harith boasted he could rival the Qur'an --- he was executed at Badr.

\textbf{Two shields against punishment.} (1) The Prophet's presence. (2) Seeking forgiveness (istighfar). After the Prophet migrated, punishment came at Badr.
\end{themebox}

\section{Theme 5: Deserved Punishment and Futile Spending (8:34--37)}
\begin{themebox}[Key Lessons]
\textbf{Obstructing the Sacred Mosque.} The Quraysh's sin was not just disbelief but preventing believers from al-Masjid al-Haram. ``Its true guardians are not but the righteous.''

\textbf{Ritual without meaning.} Their ``prayer'' at the Ka`bah was whistling and handclapping --- empty ritual devoid of devotion.

\textbf{Wealth against truth becomes regret.} The Quraysh spent vast wealth after Badr --- only to be defeated again. Allah will separate the wicked from the good and heap them into Hell.
\end{themebox}

\section{Cross-Thematic Reflections}
\begin{themebox}[Connections]
The passage moves from \textbf{exhortation} (24--25) to \textbf{remembrance} (26--28), to \textbf{promise} (29), to \textbf{historical narrative} (30--33), and finally to \textbf{prophetic declaration} (34--37). The pivot is verse 8:29: \textit{taqwa yields furqan}. This is the criterion by which the two paths are distinguished --- the very principle that verse 8:37 enacts cosmically.
\end{themebox}

\chapter{References and Sources}

\section{Primary Sources}
\begin{enumerate}
\item \textbf{The Holy Qur'an} --- Surah Al-Anfal (8), verses 24--37. Uthmani script.
Source: Al Quran Cloud API (\texttt{quran-uthmani-quran-academy}), cross-referenced with Quran.com (\texttt{text\_uthmani}).

\item \textbf{Tafsir Ibn Kathir (Abridged)} --- Imam Isma'il bin `Umar Ibn Kathir (d. 774 AH / 1373 CE). Source: Al Quran Cloud API.

\item \textbf{Tafsir Muyassar} --- King Fahd Complex for the Printing of the Holy Qur'an. Source: Al Quran Cloud API.
\end{enumerate}

\section{Translations Consulted}
\begin{enumerate}
\item Sahih International (\texttt{en.sahih}) --- English, primary
\item Marmaduke Pickthall (\texttt{en.pickthall}) --- English
\item Abdullah Yusuf Ali (\texttt{en.yusufali}) --- English
\item Muhammad Asad (\texttt{en.asad}) --- English
\item Abdullah Muhammad Basmeih (\texttt{ms.basmeih}) --- Malay
\end{enumerate}

\section{Hadith Collections Cited in Tafsir}
Sahih al-Bukhari (d. 256 AH) | Sahih Muslim (d. 261 AH) | Musnad Ahmad (d. 241 AH) | Sunan at-Tirmidhi (d. 279 AH) | Sunan an-Nasa'i (d. 303 AH) | Sunan Abu Dawud (d. 275 AH) | Sunan Ibn Majah (d. 273 AH)

\section{Corpus Provenance}
All data extracted from \textbf{Quran Wiki raw corpus}: \texttt{/root/tarbiyyah/quran-wiki/}
\begin{itemize}
\item Arabic: \texttt{raw/arabic/008-024.md} -- \texttt{raw/arabic/008-037.md}
\item Translations: \texttt{raw/translations/surah\_008/ayah\_024.json} -- \texttt{ayah\_037.json}
\item Tafsir: \texttt{raw/tafsir/008.jsonl} (lines 24--37)
\end{itemize}
Built 2026-05-10 from verified API sources. No live API calls during this research.

\end{document}
